\documentclass[11pt,letterpaper]{article}

%% =====================================================================
%%  PopRetrieve -- submission version
%%
%%  Placeholders that must be filled before submission are marked
%%  \TODO{...} and render as [[ ... ]] in the PDF so none can be missed.
%%  Set \submissionfalse to hide line numbers for a clean reading copy.
%% =====================================================================

%% ---- Packages ----
\usepackage[utf8]{inputenc}
\usepackage[T1]{fontenc}
\usepackage{mathpazo}
%% Side margins unchanged (line length is unchanged); top and bottom are tightened so the
%% full-width figure floats fit the text block. The composites are built at a fixed physical
%% width and their type is verified at that size, so the space has to come from page height
%% rather than from scaling the figures down.
\usepackage[left=2cm,right=2cm,top=1.3cm,bottom=1.3cm]{geometry}
\usepackage{graphicx}
\usepackage{float}
\usepackage{caption}
% Captions: single-spaced and one size down, as Nature sets them. The double-spacing the
% submission switch applies to body text is for reviewers to write between lines of PROSE; applying
% it to a 600-word figure legend pushes the float past a full page. singlelinecheck=false skips
% TeX's single-line measuring box, which overflows on captions this long.
\captionsetup{singlelinecheck=false, font={footnotesize, stretch=0.95}, skip=4pt, belowskip=0pt}
%% Float-page geometry. Every main figure is a single float holding the graphic and its legend,
%% so the pair has to fit \textheight. Top-aligning the float-page glue also stops LaTeX
%% spreading a lone float down the page, so any slack collects at the bottom as a page tail
%% instead of as a band of white between the graphic and its legend.
\makeatletter
\setlength{\@fptop}{0pt}
\setlength{\@fpsep}{12pt}
\setlength{\@fpbot}{0pt plus 1fil}
\makeatother
\usepackage{amsmath,amssymb}
\usepackage{booktabs}
\usepackage[dvipsnames]{xcolor}
\usepackage[most]{tcolorbox}
\usepackage{hyperref}
\usepackage[numbers,sort&compress,super]{natbib}
\usepackage{enumitem}
\usepackage{lineno}
\usepackage{etoolbox}

%% ---- Submission switches ----
\newif\ifsubmission
\submissiontrue          % line numbers + double spacing for review
\ifsubmission
  \linenumbers
  \usepackage{setspace}
  \onehalfspacing
\fi

%% ---- Unfilled-content marker ----
\newcommand{\TODO}[1]{{\color{BrickRed}\textbf{[[#1]]}}}

%% \CONTROLRANGE was removed here on 2026-09-06. Its comment claimed it was the single source of
%% truth for three quotations of the control-drug range, but it was expanded zero times in this
%% file: the range is quoted only in Supplementary Note 1, which carries its own mirrored
%% definition and now reads through it.

%% Tahoe-100M plate 3 pilot. Every value below is written by
%% analysis/tahoe_pilot/tahoe_summary_numbers.py from the per-condition and per-pair CSVs in
%% results/tahoe_pilot/, and is quoted through a macro so the Results text, the Fig. 5h-k
%% caption entries and the Methods cannot drift apart. Re-run that script to recheck any of them.
\newcommand{\TAHOELINES}{45}
\newcommand{\TAHOEDRUGS}{92}
\newcommand{\TAHOECONDS}{3630}
\newcommand{\TAHOEGONE}{0.739}
\newcommand{\TAHOEGONEQ}{0.600 to 0.836}
\newcommand{\TAHOEBELOWREAL}{0.9}
\newcommand{\TAHOEBELOWMIX}{0.11}
\newcommand{\TAHOEGTWOCEIL}{0.853}
\newcommand{\TAHOEGTWOUNSUP}{0.611}
\newcommand{\TAHOEGTWOGAP}{0.157}
\newcommand{\TAHOEGTWOPAIRS}{960}
\newcommand{\TAHOEGTWOLINES}{48}
\newcommand{\TAHOEGTWOLINESPOS}{48}
\newcommand{\TAHOEKMEANS}{0.592}
\newcommand{\TAHOEGMM}{0.582}
\newcommand{\TAHOELEIDEN}{0.563}
\newcommand{\TAHOEGTHREECC}{0.841}
\newcommand{\TAHOEGTHREECCN}{44}
\newcommand{\TAHOEGTHREEST}{0.778}
\newcommand{\TAHOEGTHREESTN}{45}
\newcommand{\TAHOEGTHREESTOPEN}{2}
\newcommand{\TAHOEGTHREESTMIN}{0.236}
\newcommand{\TAHOEGTHREEOVERLAP}{0.905}
\newcommand{\TAHOECOUPLING}{0.620}
\newcommand{\TAHOECOUPLINGN}{44}
\newcommand{\TAHOEOPENGONE}{0.394 to 0.534}
\newcommand{\TAHOEOPENGTHREE}{0.605 to 0.764}
\newcommand{\TAHOECOUPLINGCROSS}{0.129}
\newcommand{\TAHOEPARTAGREE}{0.466}

%% The premise correlation, in two forms (analysis/natural/zhao_premise_disjoint.py;
%% results/zhao_gbm/premise_disjoint.csv). The first correlates the mean signature with the
%% malignant-compartment response, and the malignant compartment is HALF of that mean signature by
%% construction, so part of it is arithmetic. The second correlates two compartments that share no
%% cells and is the load-bearing figure. Since 2026-09-06 the Results quote the disjoint one only,
%% which is what Fig. 4f draws and what its caption has always described, and the overlapping form
%% sits in the Methods beside the design it motivates.
%% \PREMISESHARE WAS RETIRED 2026-09-06 AND MUST NOT BE REINSTATED AT 43%%. It read "the malignant
%% cells are 43%% of the cells entering that mean signature", which is wrong twice over. The mean
%% signature is 0.5*(malignant delta + myeloid delta), an equal-weight average of two compartment
%% deltas rather than a cell-weighted mean (analysis/natural/zhao_two_gates.py:129), so malignant
%% contributes exactly HALF of it by construction. And 43%% was 41,314/96,225, the malignant share
%% of all called cells across THREE compartments, a denominator holding 23,233 oligodendrocytes
%% that never enter this signature. By cell count the malignant share of the drug-treated cells of
%% the four patients in this analysis is 60.2%%, not 43%%. The overlap is now stated in the Methods
%% as the arithmetic half-weight, which is what the code does.
\newcommand{\PREMISEOVERLAP}{+0.878}
\newcommand{\PREMISEDISJOINT}{+0.835}

%% Class C, semantically matched functional oracle (analysis/class_c/class_c_functional_oracle.py;
%% results/upgrade/class_c_functional_oracle.json). Defined once so no number can drift between
%% the text, the figure caption and the Methods.
\newcommand{\CTXLINES}{966}
\newcommand{\NGDSCDRUGS}{286}
\newcommand{\NCQUERIES}{103}
%% ---- Phase-II intervention task, read from results/phase2_transition/eligibility/task_freeze.json
%% Added 2026-09-04 with the Methods subsection. These four counts appear in the Abstract, the
%% Introduction, two Results paragraphs, the Fig. 5 and Fig. 6 captions and the Methods; every one
%% of those was a hand-typed literal, which is the failure mode recorded in A5d and A5f.
\newcommand{\NCAND}{92}                 % candidate compounds in the fixed library
\newcommand{\NQUERIES}{3,992}           % eligible queries
\newcommand{\NCTX}{44}                  % usable held-out cell lines
\newcommand{\NFITCTX}{43}               % training contexts per held-out line
\newcommand{\NSEEDS}{five}              % fixed seeds; 19,960 query--seed pairs
\newcommand{\NQSEED}{19,960}
%% THE THREE CLASS-C VALUES BELOW THAT TOUCH AN ENERGY KERNEL WERE REISSUED 2026-09-04, and they
%% are the last of the missed U-statistic reissues. results/_pre_ustat_backup vs results/upgrade
%% of class_c_functional_oracle.csv differ in EXACTLY three columns: energy_rho (+0.27609 to
%% +0.26504), energy_rho_partial_magmatch (+0.09693 to +0.10528) and energy_rho_uncentered
%% (+0.09672 to +0.09024). The other four Class-C rankings compare mean vectors or scalars, never
%% a population, so they cannot move and did not. Figure 3h had already been rebuilt and drew
%% +0.265 and +0.105; these macros and Supplementary Table 3 kept the V values, which is how the
%% Results paragraph and its own figure came to print different numbers for the same panel.
\newcommand{\ENERGYFUNC}{+0.265}
\newcommand{\MEANCTRLFUNC}{+0.083}
\newcommand{\MEANRAWFUNC}{+0.241}
\newcommand{\MAGMATCHFUNC}{+0.232}
\newcommand{\POTMATCHFUNC}{+0.399}
\newcommand{\ENERGYPARTIAL}{+0.105}
\newcommand{\ENERGYWINS}{62}
\newcommand{\ENERGYWINSP}{0.054}
\newcommand{\CENTERUNCEN}{85\% of drug pairs positive, median $+0.178$, versus 52\% and
  $+0.015$ once corrected}

%% ---- Partial-observed retrieval: denominators and headline statistics ----
%% Every value below is verified against results/exp12_partial_observed_retrieval/per_query_scores.csv
%% by analysis/audit/numerical_contract_stage{1,2}.py; the audit is
%% manuscript/reference/numerical_contract_audit.md. THE FOUR DENOMINATORS ARE TWO INDEPENDENT
%% CUTS OF THE SAME 765 QUERIES, NOT A NESTED HIERARCHY:
%%   by evaluator definability : 600 MoA-defined (split_type == leave_drug_out) + 165 undefined
%%   by gate label             : 621 recommended + 133 non-recommended + 11 mean-sufficient
%% A value computed on one cut must never be printed beside a value computed on the other. The
%% published range "+0.056 to +0.129" did exactly that (lower bound from the 621, upper from the
%% 600) and is corrected here to \CLASSARANGE, wholly on the 600.
\newcommand{\NALL}{765}                 % all partial-observed queries; sha256 11e6c902...
\newcommand{\NMOA}{600}                 % Class B defined; sha256 1dda00b1...
\newcommand{\NMOAEXCL}{165}             % 120 leave_MoA_out + 45 partial_library
\newcommand{\NREC}{627}                 % gate-recommended (was 621 under the V arm)
\newcommand{\NNONREC}{127}              % non-recommended (mean_or_no_call); was 133
\newcommand{\NMEANSUFF}{11}             % mean-sufficient

%% Class A, worst-case coverage against mean cosine. \REGRETALL is on \NALL, the rest on \NMOA.
%% REISSUED 2026-09-03 under the unbiased U-statistic energy distance. The V-statistic values
%% were +0.118 / +0.258 / 72\% / 1.6e-66; see docs/phase2/POST_REPAIR_MASTER_RESULTS.md A5.
\newcommand{\REGRETALL}{+0.031}         % median, n = 765
\newcommand{\REGRETALLCI}{+0.017\text{ to }+0.049}   % 95\% bootstrap CI of that median
\newcommand{\REGRETALLMEAN}{+0.132}     % mean,   n = 765
\newcommand{\REGRETALLFRAC}{56\%}       % improved, n = 765
\newcommand{\REGRETALLP}{1.6\times10^{-12}}
%% REISSUED 2026-09-03 under the U arm, in the same pass as \REGRETALL above. They were NOT
%% reissued with it: the comment above covered only the \REGRETALL group, and these eight kept
%% their V-statistic values while the Results prose carried hand-typed copies of them. The
%% Class-A median is the one that moved far, +0.129 to +0.028, because coverage-worst is the
%% scorer the estimator moved most (Fig. 2a: Hit@1 0.589 to 0.747). All eight are recomputed from
%% results/exp12_partial_observed_retrieval/per_query_scores.csv on split_type == leave_drug_out.
\newcommand{\CLASSAMED}{+0.028}         % median, n = 600 (V: +0.129)
\newcommand{\CLASSAMEAN}{+0.135}        % mean,   n = 600 (V: +0.275)
\newcommand{\CLASSBMED}{0.000}          % median, n = 600 (V: 0.000, unchanged)
\newcommand{\CLASSBMEAN}{-0.024}        % mean,   n = 600 (V: -0.037)
\newcommand{\CLASSBWORSE}{40.8\%}       % V: 41.3\%
\newcommand{\CLASSBBETTER}{35.7\%}      % V: 34.2\%
\newcommand{\CLASSBTIED}{23.5\%}        % V: 24.5\%
\newcommand{\CLASSBP}{7.7\times10^{-3}} % two-sided Wilcoxon, zeros dropped, 459 non-zero of 600.
%% Used in BOTH the Results sentence and the Fig. 3a caption entry, which used to carry two
%% hand-typed copies that disagreed. Under the V arm it was 2.5e-4.
\newcommand{\CLASSARANGE}{$+0.028$ to $+0.067$}   % five variants, MEDIAN, n = 600 (V: +0.053 to +0.129)
\newcommand{\CLASSBRANGE}{$-0.011$ to $-0.033$}   % five variants, MEAN, n = 600, same queries
%% A and B are summarised by different statistics on purpose: all five Class-B MEDIANS are
%% exactly 0.000, so a median range would print "0.000 to 0.000" and say nothing about the five.
%% Rank correlation between a variant's Class-A median gain and its Class-B mean gain, across the
%% five variants on the same 600 queries. n = 5 is far too small to test, so this is reported as a
%% description and never with a p-value. It was $-0.90$ under the V arm and is $+0.60$ under U,
%% which is why no sentence in the paper rests on it: at n = 5 this quantity is not stable enough
%% to carry an argument, and it carries none.
\newcommand{\VARIANTCOUPLING}{$+0.60$}
\newcommand{\GATEREC}{+0.032}           % median, n = 627
\newcommand{\GATENONREC}{+0.040}        % median, n = 127
%% The gate's null, measured rather than eyeballed (figures/fig2/fig2d.py computes all three
%% from results/exp12_partial_observed_retrieval/per_query_scores.csv with a seeded bootstrap,
%% so the panel and this preamble cannot disagree). Direction: recommended MINUS non-recommended.
\newcommand{\GATEDIFF}{-0.008}
\newcommand{\GATEDIFFCI}{$-0.032$ to $+0.056$}
\newcommand{\GATEDIFFP}{0.64}

%% Real-data projection (results/exp13_real_data_projection/projection.csv).
%% \PROJP was published as 1.3e-7, which is the ONE-sided half of the two-sided value; every
%% other paired test in this paper is two-sided, so the two-sided value is used.
\newcommand{\PROJMEAN}{+0.0011}
\newcommand{\PROJP}{2.6\times10^{-7}}
\newcommand{\PROJN}{239}
\newcommand{\PROJTIED}{102}
\newcommand{\PROJNONZERO}{137}

%% Class C, GDSC2 functional oracle. \ENERGYPARTIAL is a PARTIAL Spearman (magnitude channel
%% removed) and \ENERGYUNCEN is a RAW Spearman under the uncentered oracle. Under the V arm they
%% rounded to the same three decimals (+0.097 and +0.0967) while being quoted in adjacent
%% paragraphs, which is why the uncentered one carries a fourth decimal. Under the U arm they no
%% longer collide (+0.105 and +0.0902); the fourth decimal and the test-type labels are kept
%% anyway, because the collision was a property of the numbers rather than of the wording.
%% Coupling between an energy distance and the candidate's OWN response magnitude. Used by R5 and
%% by the SI confounder audit; previously a bare literal in both.
%% CORRECTED 2026-09-04 from +0.791. The Fig. 3k caption printed +0.780 for the same panel and
%% the same quantity; figures/source_data/fig3hi_class_c_potency.csv gives a median
%% energydist_vs_candmag_rho of +0.7797 over the 103 queries, so the caption was right and this
%% macro was carrying a value nothing in the repository produces. The caption entry now reads
%% through this macro rather than through a second hand-typed copy.
\newcommand{\ENERGYCANDMAG}{+0.780}
\newcommand{\ENERGYUNCEN}{+0.0902}
\newcommand{\MEANCTRLUNCEN}{+0.138}
\newcommand{\MAGMATCHUNCEN}{+0.217}

%% Oracle-shape test (results/upgrade/oracle_shape_test.json).
\newcommand{\SHAPEN}{659}
\newcommand{\SHAPENCTRL}{213}
\newcommand{\SHAPENIFN}{227}
\newcommand{\SHAPENCO}{219}
%% REISSUED 2026-09-03 under the unbiased U-statistic. The V-arm values were
%% +0.242 / +0.146 / +0.529 / +0.334 / +0.125 / +0.352 / +0.177 / +0.022. Two of the eight do not
%% move at all: neither the mean cosine nor the magnitude scalar passes through an energy
%% distance, so their scores under the MEAN-shaped evaluator are identical across arms. See
%% docs/phase2/FINAL_FIGURE_NUMBER_LEDGER.md, Figure 4 panel b.
\newcommand{\SHAPEMEANORACLEMEAN}{+0.242}
\newcommand{\SHAPEMEANORACLEENERGY}{+0.222}
\newcommand{\SHAPEDISTORACLEENERGY}{+0.393}
\newcommand{\SHAPEDISTORACLEMEAN}{+0.230}
\newcommand{\SHAPEMAGMEAN}{+0.125}
\newcommand{\SHAPEMAGDIST}{+0.276}
\newcommand{\SHAPEEDGEDIST}{+0.117}
\newcommand{\SHAPEEDGEMEAN}{+0.098}
%% Frangieh protein panel and join. TODO(U-2) RESOLVED 2026-09-03. The barcode-join count is now
%% persisted to results/upgrade/oracle_shape_join_provenance.json and Figure 4 panel a reads it,
%% asserting that the three counts agree: the protein file, the RNA file and their intersection
%% are all 218,331, so the analysis unit is not smaller than either input and the published
%% figure stands. Supplementary Table 2's 218,023 is the same RNA h5ad under one further filter,
%% not a different source: 218,331 - 308 = 218,023, and the 308 cells are the nine knockout targets
%% dropped by --min-cells 50 (248 KOs -> 239). CORRECTED 2026-09-06: this note used to attribute the
%% gap to the 2,000-HVG step, which reduces genes (23,712 -> 2,000) and never removes a cell.
%% Both paths run the identical normalize_total(1e4)/log1p/HVG(2000, seurat) chain on the same input.
\newcommand{\SHAPECELLS}{218{,}331}
\newcommand{\PROTMARKERS}{20}           % 24 measured, 4 isotype controls dropped

%% ---- R3 Hit@1 values (results/exp08_signature_baselines/summary.csv) ----
%% Unweighted macro-mean over the seven (task x setting) cells is PRIMARY; the query-weighted
%% variant is a sensitivity analysis. \HITMEAN serves both mean_cosine and cmap_cosine because
%% they are the same number to six figures (0.388492); one macro makes that identity unbreakable
%% by editing. The six-figure form is quoted only in R2, where the identity is itself the result.
%% exp08 Hit@1, macro-mean over the seven task--setting combinations unless marked W (weighted by
%% query count). REISSUED 2026-09-03 under the unbiased U-statistic: \HITENERGY was 0.837 and
%% \HITENERGYW 0.887, both V-statistic values. The other seven do not touch an energy distance and
%% cannot move. All nine were defined and unreferenced while the Results prose hand-typed copies of
%% them, which is why the two stale ones survived the reissue; the prose now reads the macros.
\newcommand{\HITENERGY}{0.836}
\newcommand{\HITMAG}{0.788}
\newcommand{\HITMEAN}{0.388}
\newcommand{\HITPCADIST}{0.778}
\newcommand{\HITPCAMEAN}{0.518}
\newcommand{\HITENERGYW}{0.884}
\newcommand{\HITMEANW}{0.421}
\newcommand{\HITFRANGMEAN}{0.600}
\newcommand{\HITFRANGMAG}{0.611}
\newcommand{\HITFRANGENERGY}{0.578}
\newcommand{\HITWTCS}{0.4635}
%% The two steps of the information ladder, differences of the three macros above.
\newcommand{\HITSTEPMAG}{+0.399}
\newcommand{\HITSTEPDIST}{+0.048}

%% ---- Section formatting ----
\usepackage{titlesec}
\titleformat{\section}{\large\bfseries}{}{0em}{}
\titleformat{\subsection}{\normalsize\bfseries}{}{0em}{}
\titlespacing*{\section}{0pt}{1.5ex plus 0.5ex}{0.8ex plus 0.2ex}
\titlespacing*{\subsection}{0pt}{1.2ex plus 0.3ex}{0.5ex plus 0.1ex}

\hypersetup{colorlinks=true,linkcolor=MidnightBlue,citecolor=MidnightBlue,urlcolor=MidnightBlue}
\setlength{\emergencystretch}{3em}

%% ---- Float placement ----
%% Without these, a figure that is tall relative to the text block gets promoted to a float page and
%% strands the rest of the sheet blank (Fig. 3 was doing exactly that). LaTeX's defaults reserve 20%
%% of every page for text (textfraction) and refuse to let a top float exceed 70% (topfraction), so a
%% 6.9 in figure plus a legend has nowhere legal to sit and is deferred. These values let a large
%% float share its page with running text instead.
\renewcommand{\topfraction}{0.95}
\renewcommand{\bottomfraction}{0.95}
\renewcommand{\textfraction}{0.06}
\renewcommand{\floatpagefraction}{0.80}
\setlength{\textfloatsep}{14pt plus 3pt minus 3pt}
\setlength{\floatsep}{12pt plus 3pt minus 3pt}
\setcounter{topnumber}{2}
\setcounter{bottomnumber}{1}
\setcounter{totalnumber}{3}

\usepackage{placeins}

\begin{document}

%% =====================================================================
%%  TITLE PAGE
%% =====================================================================
\thispagestyle{empty}
\begin{flushleft}

{\LARGE\bfseries When do single-cell response distributions improve drug retrieval? \par}


\end{flushleft}
\bigskip
\noindent\rule{\textwidth}{0.4pt}


%% =====================================================================
%%  ABSTRACT
%% =====================================================================
\begin{abstract}
\noindent

Artificial intelligence virtual cells aim to predict cellular responses to perturbations and, ultimately, to guide intervention selection. Conventional transcriptomic drug retrieval represents each perturbation by a single expression signature, whereas single-cell perturbation atlases preserve cell-to-cell variation in perturbation responses. Here we define population-to-population drug retrieval (PopRetrieve), in which candidate drugs are ranked by comparing single-cell response populations, and ask when this additional information improves drug ranking. We find that most of the gain over direction-only mean matching is recovered by a mean-based score that also retains response magnitude, with a smaller additional gain from retaining the full response population. The advantage of population-level retrieval weakens as evaluation moves away from response matching: it disappears or reverses under mechanism-of-action recovery, and most of its association with an external functional readout is reproduced by response magnitude alone. The same held-out protein measurements favour different retrieval strategies depending on whether they are summarized as means or retained as populations, showing that the statistical form of an external evaluator can influence which retrieval strategy appears preferable. With candidate responses observed, population-level retrieval corrects an incorrect top-ranked choice about four times as often as it overturns a correct one relative to magnitude-aware mean matching, and about five times relative to direction-only matching, but the gain is confined to the roughly one-fifth of queries for which mean matching leaves room for improvement. When candidate responses are predicted rather than observed, this oracle advantage disappears and can reverse sign even when mean-level retrieval performance is preserved. Together, these results show that the value of single-cell population information depends on whether information beyond the mean changes candidate ordering and, when required, survives forward prediction.

\end{abstract}

Keywords: artificial intelligence virtual cell; single-cell perturbation; drug retrieval; population-to-population retrieval; cellular heterogeneity; objective--utility mismatch


\bigskip\noindent\rule{\textwidth}{0.4pt}\bigskip

%% =====================================================================
%%  INTRODUCTION
%% =====================================================================
\section*{Introduction}

Artificial intelligence virtual cells\cite{AIVirtualCell,VirtualCellChallenge} aim to predict cellular responses to perturbations\cite{ref7,ref8,ref9} and, ultimately, to help prioritize interventions for testing. These goals pose two related problems. Forward prediction asks what response a perturbation will produce\cite{ref4,ref5,ref6,ref10}; retrieval asks which perturbation best matches a cellular response of interest\cite{ref12}. Transcriptomic drug retrieval has traditionally represented each drug by a single perturbation signature and ranked candidates by similarity or reversal\cite{ref1,ref2}. This approach has made large drug libraries searchable, but collapses variation across cells into a single summary. Single-cell perturbation atlases now capture this variation\cite{ref3,dataSciPlex3,Tahoe100M}, including state-dependent responses\cite{dataZhao} and effects confined to minority cell populations\cite{RareCellVariability}. Here we define population-to-population drug retrieval (PopRetrieve), in which candidate drugs are ranked by comparing single-cell response populations, and ask when this additional information improves drug ranking.

More detailed response information does not necessarily lead to a better drug ranking. Two drugs can have similar average responses yet differ in how individual cells respond, creating information that a population-aware method can use. But an apparent advantage over classical signature matching can arise for more than one reason. Population comparison retains response magnitude that directional similarity discards, so a gain over direction-only matching need not arise from within-population structure. The apparent benefit also depends on how retrieval quality is evaluated: a criterion that preserves full response distributions\cite{OptimalMetrics,StandardizedEval} will naturally favour information that a population-aware scorer retains\cite{CellOT,CellFlow,Departures}. Such gains establish better response matching, but not necessarily a better biological or functional ranking\cite{ref13,ref14,ref15,ref16,MetricGaming}.

We therefore decompose the problem into information retained, information that changes candidate ordering, and information that survives prediction. We first distinguish three scoring rules by how much of the response each retains: direction, direction plus magnitude, and the full response population. This separates gains recovered by restoring response magnitude from those that require information beyond the mean. We then evaluate the resulting rankings using increasingly task-distal criteria: response matching, mechanism recovery and external functional readouts\cite{GDSC,GDSCLandscape}. We use state-dependent perturbation responses as a concrete mechanism through which population information can affect retrieval. For such information to matter, it must be present and recoverable, change candidate ordering beyond what the mean already determines and, when candidate responses are not observed, survive forward prediction. This separates the existence of single-cell response structure from its value for intervention selection.

Here we show that most of the advantage of population comparison over direction-only signature matching is recovered by a mean-based score that also retains response magnitude, leaving a smaller gain that requires the full response population. The advantage of population comparison weakens as evaluation moves away from response matching: it disappears or reverses under mechanism-of-action recovery, and most of its association with an external functional readout is reproduced by response magnitude alone. The same held-out protein measurements favour different retrieval strategies depending only on whether they are summarized as means or retained as populations, showing that external evaluation is not automatically neutral. In a source-to-target intervention-identification task, population information provides a smaller additional gain beyond magnitude-aware mean matching when candidate responses are observed, and its whole advantage over direction-only mean matching is confined to the queries for which that baseline leaves room for improvement. That advantage does not survive when candidate responses must instead be predicted. Together, these results show that the value of single-cell population information depends not simply on whether additional cellular structure exists, but on whether it changes candidate ordering and survives to the final retrieval decision.


%% =====================================================================
%%  RESULTS
%% =====================================================================
%% FIGURE PLACEMENT, 2026-09-05. Each main figure is ONE [htbp] float holding its graphic and
%% its legend, declared immediately after the paragraph that first cites it. The earlier
%% arrangement, a bare \includegraphics preceded by \clearpage with a hand-numbered legend
%% paragraph after it, is what produced the blank space: the \clearpage left the preceding text
%% page 20 to 40%% full, and a legend too long to follow its graphic spilled onto the next page
%% (p. 17 of the 2026-09-05 build carried one legend line and nothing else). A single float can
%% hold both only while the legend fits what the graphic leaves of \textheight, so the legends
%% are held to one sentence per panel and the page height was raised in the geometry above.
\section*{Results}

\subsection*{Population-to-population retrieval generalizes classical signature matching}

Population-to-population retrieval ranks candidates by comparing query and candidate response populations directly (Fig.~\ref{fig:1}a,b), whereas classical mean-signature retrieval\cite{ref1,ref2} first collapses each population to a single mean differential-expression vector. Retaining the cells makes response structure available that averaging removes, including differences between cellular states or subpopulations. The retrieval problem is unchanged; what changes is how much of the measured response reaches the ranking.

Two candidate populations can share the same mean response and still differ markedly in how individual cells respond (Fig.~\ref{fig:1}c,d). A population-aware score separates such candidates where a score computed from the mean cannot, because any such score reads only the mean the two candidates share. We therefore hold the candidate rankings fixed and judge them under criteria that progressively reduce their coupling to the retrieval score: response matching, mechanism recovery and external functional readouts (Fig.~\ref{fig:1}e,f).

% FIGURE 1
\begin{figure}[htbp]
\centering
\includegraphics[width=\textwidth]{figures/fig1.pdf}
\caption{\textbf{Perturbation responses can be represented and evaluated at multiple information resolutions.}
\textbf{a}--\textbf{f}, schematic; \textbf{g}, a seeded synthetic construction; \textbf{h}, an analytic illustration.
\textbf{a}, Three scoring rules read one single-cell response population (grey, two states). Glyphs, the geometry each rule uses: a mean vector $\mu$ normalised to a unit circle, the same vector at its true length, and the population $P$ with $\mu$ still on it (orange diamond, mean; blue cells, population). Matrix, the properties each rule retains (filled) or discards (open); the columns name properties and carry no quantity.
\textbf{b}, One pair of populations and the three scoring rules: the mean route keeps one vector each, the population route compares the cells.
\textbf{c}, Under drug B a resistant fifth does not move, yet the two treated means coincide in direction and in magnitude; faint cells, untreated.
\textbf{d}, Three populations sharing one sample mean exactly, so both mean scores tie by construction while the population score, an energy distance over the drawn cells, prefers A.
\textbf{e}, Objective-aligned, the score that produced the ranking also grades it; less score-aligned, the same ranking is judged by mechanism of action, target engagement or GDSC2 viability (icons).
\textbf{f}, Three evaluation levels ordered by coupling to the retrieval score; filled, that level is a study's headline metric.
\textbf{g}, Energy distance against the mean-only limit as the spread $\lambda$ is restored, closing at $\lambda=0$, where the energy distance between two point masses is exactly twice the Euclidean distance between their means.
\textbf{h}, Four per-subpopulation distances 0.20, 0.55, 0.80 and 1.30, the worst-matched state crossed; frames, the weight the aggregate places on each state at $\beta\to0$, $\beta=2$ and $\beta\to\infty$.}
\label{fig:1}
\end{figure}


Mean-signature and population-level retrieval are one task at two resolutions, connected by the amount of within-population variation retained. As that variation vanishes, a response population collapses onto its mean, and the energy distance between two such populations differs from the Euclidean distance between their means only by a factor of two, which cannot reorder candidates. The magnitude-aware mean L2 is therefore the zero-variance limit of energy-based population retrieval (Fig.~\ref{fig:1}g; Methods). The cosine used by classical connectivity scoring discards the length of that mean difference and coarsens the same representation one step further (Fig.~\ref{fig:1}a,b). The implemented CMap-style cosine applies that same rule to the same mean vectors and is therefore identical to mean-cosine retrieval, whereas the canonical L1000 weighted connectivity score remains a distinct rank-based score (Fig.~\ref{fig:2}f).

The coverage family provides a second continuum: sharpening its aggregation moves the score from averaging subpopulation distances toward the worst-matched state, and a single undivided component recovers the global distance (Fig.~\ref{fig:1}h; Methods). These relationships place classical signature matching and population comparison on one axis of sensitivity to within-population structure (Supplementary Table~1).


\subsection*{Most of the gain over directional matching does not require the full response population}

Population-level retrieval substantially outperformed direction-only mean matching, and most of that gain did not require the full response population. With candidate responses directly observed, so that the full cellular populations are available to both approaches, energy-distance retrieval reached a Hit@1 of \HITENERGY, macro-averaged across the seven task--setting combinations analysed here, against \HITMEAN\ for mean-cosine retrieval and the numerically identical CMap-style cosine baseline\cite{ref1,ref2}. Between them sits a magnitude-aware mean L2 at \HITMAG, applied to the same mean signatures: the step to that score is $\HITSTEPMAG$ of the gap and the step from it to the full distribution a further $\HITSTEPDIST$ (Fig.~\ref{fig:2}a). Because that first step changes both what the score reads, from direction alone to direction and magnitude, and the geometry it uses, from angular similarity to Euclidean distance, mean L2 is a magnitude-aware mean baseline rather than a pure magnitude ablation. The same ordering held in PCA space (Fig.~\ref{fig:2}a).

% FIGURE 2
\begin{figure}[htbp]
\centering
\includegraphics[width=\textwidth]{figures/fig2.pdf}
\caption{\textbf{A magnitude-aware mean score recovers most of the gain beyond directional scoring.}
Points and intervals are defined in Methods.
\textbf{a}, Hit@1 for nine scores, macro-averaged over seven task--setting combinations: direction-only mean cosine \HITMEAN, magnitude-aware mean L2 \HITMAG\ and energy distance \HITENERGY, so the two steps are $\HITSTEPMAG$ and $\HITSTEPDIST$. Dagger, the published CMap connectivity score (WTCS).
\textbf{b}, The same three scores per task ($n=90$, 1,080 and 90 queries); the magnitude step is $+0.52$ and $+0.41$ on the two SciPlex3 tasks and $+0.01$ on Frangieh, where the distributional step is $-0.03$.
\textbf{c}, Worst-case subpopulation coverage against direction-only mean cosine over \NALL\ partial-observation queries; curve, the quantile function of the per-query regret reduction, positive favouring coverage (less regret on $\REGRETALLFRAC$ of queries, more on 36\%, ties on 8\%). Median $\REGRETALL$ (band, 95\% bootstrap CI $\REGRETALLCI$; paired Wilcoxon $P=\REGRETALLP$, mean $\REGRETALLMEAN$), and $+0.035$ on the 165 queries with a genuinely incomplete library ($P=6.0\times10^{-5}$).
\textbf{d}, Energy advantage over mean cosine against the minority fraction $\alpha$ ($n=20$ seeds per point; top, query composition), falling from $+1.00$ to $0.00$ in K562 and not falling in A549 or MCF7 ($+0.40$ to $+0.65$; $+0.35$ to $+0.45$).
\textbf{e}, Median regret reduction for five population-level scores on \NREC\ diagnostic-positive queries. Against direction-only mean cosine all five 95\% bootstrap intervals exclude zero ($\GATEREC$ to $+0.067$); against the magnitude-aware mean L2 the three distributional distances have a median of exactly zero and both coverage scores are negative ($-0.030$, $-0.038$).
\textbf{f}, Spearman correlation among seven scores over 54,180 query--candidate scores from 1,260 queries, an agreement between scoring rules rather than retrieval evidence; shade, $|\rho|$; rings, the closest pair within a tier. The magnitude-aware mean L2 sits with the distributional scores ($\rho=0.68$ to $0.80$), not with the direction-only scores ($\rho=0.17$ to $0.20$).}
\label{fig:2}
\end{figure}

This ordering was not universal. On Frangieh Perturb-CITE-seq\cite{dataFrangieh}, the only one of the seven combinations not built from SciPlex3, the magnitude-aware mean L2 was the best of the nine scores (Hit@1 \HITFRANGMAG), ahead of mean cosine (\HITFRANGMEAN) and of energy distance (\HITFRANGENERGY; Fig.~\ref{fig:2}b). There, retaining magnitude helped and retaining the full response distribution did not.

Population-level scoring also improved response matching under partial observation, where only part of each candidate population is available. The same rankings were judged by a regret criterion that penalizes selecting a candidate whose per-cell outcomes are worse than the best available (Methods). Across \NALL\ queries, the worst-case subpopulation-coverage score reduced median regret by $\REGRETALL$ relative to direction-only mean cosine (95\% bootstrap CI $\REGRETALLCI$), with lower regret on $\REGRETALLFRAC$ of queries (Fig.~\ref{fig:2}c).

Reducing the minority response state from half of the query population to a tenth did not erode the energy-distance advantage consistently: it collapsed in one cell line and held in the other two (Fig.~\ref{fig:2}d). The pattern belongs to the population-aware family rather than to one score. All five variants improved on direction-only mean cosine, whereas against the magnitude-aware mean L2 the three distributional distances have a median difference of exactly zero and both coverage scores are negative (Fig.~\ref{fig:2}e).

A magnitude-aware mean score therefore recovers most of the gain beyond directional scoring, and the further increment from retaining the full population is real and much smaller. Both criteria used so far reward correspondence between response populations, so we next ask whether that increment survives biological and functional evaluation.

\subsection*{Population-level gains weaken under biological and functional evaluation}

Whether that additional information improves the same candidate rankings depends on how success is defined. We therefore kept the rankings fixed and changed only the criterion used to judge them (Fig.~\ref{fig:3}).

Mechanism-of-action recovery is the first task-proximal biological criterion. Among the \NMOA\ partial-observation queries carrying both response-matching and mechanism annotations, the worst-case subpopulation-coverage scorer improved on direction-only mean cosine by a median of $\CLASSAMED$ under response matching, whereas the identical rankings judged by mechanism-of-action recovery gave a median of $\CLASSBMED$ and a mean of $\CLASSBMEAN$ (Fig.~\ref{fig:3}a,b). At the query level, mechanism recovery decreased more often than it increased (\CLASSBWORSE\ against \CLASSBBETTER). The same divergence between the two criteria appeared in all five population-level scoring variants, whose response-matching median gains span \CLASSARANGE\ while their mechanism-recovery means span \CLASSBRANGE. Gains under response matching therefore did not consistently carry over to the recovery of biologically related drugs.

% FIGURE 3
\begin{figure}[htbp]
\centering
\includegraphics[trim=0 2 0 3.5, clip, width=\textwidth]{figures/fig3.pdf}
\caption{\textbf{The population advantage weakens, and can change direction, once the judge stops sharing the retrieval objective.}
Summary points, medians, except in \textbf{d} and \textbf{e}, where they are quartile means; bars, 95\% bootstrap intervals, except in \textbf{a}, \textbf{g} and \textbf{j}, where they are interquartile ranges; cell lines, line style in \textbf{b}, marker shape in \textbf{i}.
\textbf{a}, \NMOA\ leave-one-drug-out queries judged twice: median $\CLASSAMED$ under response matching against \CLASSBMED\ under mechanism-of-action recovery (24\% ties; Wilcoxon $p=\CLASSBP$).
\textbf{b}, Per cell line, medians exactly 0.000 in all three ($n=200$ each).
\textbf{c}, Minority-state coverage gain over 239 tasks, 102 of which score identically; dots, individual tasks; diamonds, dataset means; daggers, constructed data.
\textbf{d}, Population-minus-mean gain by divergence quartile, objective-aligned judge ($n=\NALL$); quartile means $-0.0035$ (CI $-0.0076$ to $+0.0001$), $+0.0038$, $+0.0054$, $+0.0051$, against medians $0.0000$, $+0.0004$, $+0.0012$, $+0.0020$ (Spearman $\rho=+0.14$, $p=1\times10^{-4}$).
\textbf{e}, The same quartiles under mechanism-of-action recovery ($n=\NMOA$): no interval on the quartile mean lies above zero, the two least divergent wholly below it ($-0.050$, $-0.039$; $q=0.002$, $0.025$, Benjamini--Hochberg). With 17 to 32\% of queries in each quartile scoring identically under this criterion, the corresponding medians are $-0.0062$, $0.0000$, $0.0000$ and $0.0000$; the mean is the summary drawn here and the median is the summary quoted in \textbf{a} and \textbf{b}.
\textbf{f}, Regret reduction, \NREC\ recommended queries against \NNONREC\ declined ($\GATEDIFF$, CI \GATEDIFFCI, Mann--Whitney $p=\GATEDIFFP$).
\textbf{g}, Below, structure reliability against divergence ($\rho=-0.19$, $P=1.5\times10^{-7}$); above, declined queries are \emph{more} divergent (1.68 against 1.66, $P=0.009$).
\textbf{h}, Ranking against GDSC2 drug--drug similarity ($n=\NCQUERIES$); filled, ranking methods; arc, energy partialled on magnitude.
\textbf{i}, Per query: energy beats the magnitude scalar on \ENERGYWINS\ queries and loses on 40 (paired $p=\ENERGYWINSP$).
\textbf{j}, Rankings against absolute GDSC2 potency: the two rankings computed without control subtraction, energy and raw mean cosine, turn negative ($-0.506$, $-0.533$); the control-subtracted mean cosine stays just above zero ($+0.105$) and candidate magnitude alone is highest ($+0.692$).
\textbf{k}, Per-query Spearman between energy distance and candidate response magnitude, median $\ENERGYCANDMAG$.}
\label{fig:3}
\end{figure}

A second task-proximal criterion tests how well population-level retrieval preserves minority-state responses. Across 239 task instances, population-level retrieval improved minority-state coverage by only $\PROJMEAN$ on average, and in 102 of them the two approaches scored exactly the same (Fig.~\ref{fig:3}c). The small positive effect was concentrated in the constructed cross-cell-line tasks ($+0.0037$), whereas the remaining settings were near zero or slightly negative. Stratifying by measured response divergence revealed a gradient only on the objective-aligned criterion: the minority-state advantage rose from below zero in the least divergent quartile to about $+0.005$ in the two most divergent (Fig.~\ref{fig:3}d), while no divergence stratum reached a mechanism-of-action gain whose interval lay above zero (Fig.~\ref{fig:3}e). Stronger population differences therefore do not by themselves produce a more biologically useful ranking, and the pre-specified diagnostic gate did not identify the queries that did benefit (Fig.~\ref{fig:3}f,g; Supplementary Note~1).

Population-level retrieval does carry information that an external functional readout can see. Using GDSC2\cite{GDSC,GDSCLandscape} dose--response profiles from cell lines excluded from the transcriptional retrieval experiment, and removing the shared cell-line sensitivity effect, energy-distance retrieval reached a Spearman correlation of $\ENERGYFUNC$ with drug--drug functional similarity, against $\MEANCTRLFUNC$ for the corresponding mean-signature baseline (Fig.~\ref{fig:3}h; Supplementary Table~3).

Much of that advantage did not require comparing cellular distributions. Ranking candidates only by how closely their overall response magnitude matched the query reached $\MAGMATCHFUNC$ from a single scalar per perturbation, and once that magnitude channel was accounted for the remaining association for energy retrieval was $\ENERGYPARTIAL$, positive in all three cell lines (Fig.~\ref{fig:3}h,i). Scored against absolute GDSC2 potency instead, energy and raw mean cosine turn negative, since matching a weak query's magnitude returns other weak drugs, whereas ranking by candidate magnitude alone turns positive (Fig.~\ref{fig:3}j). The channel is also directly measurable: across queries the query--candidate energy distance tracks the candidate's own response magnitude at a median Spearman correlation of $\ENERGYCANDMAG$ (Fig.~\ref{fig:3}k). Population-level retrieval therefore carries external functional signal, and most of the signal detected here is recoverable without any cell-to-cell distributional structure.

The apparent value of population-level retrieval therefore depends on how success is measured, and it does so even when the criterion moves beyond transcriptomic response matching.

\subsection*{Evaluator form can reverse the preferred retrieval strategy, and natural heterogeneity leaves little reordering headroom}

We next asked whether the dependence on the evaluation criterion persists when rankings are judged by a biological measurement that no retrieval method can see.

Frangieh Perturb-CITE-seq\cite{dataFrangieh} measures RNA and \PROTMARKERS\ surface proteins in the same \SHAPECELLS{} cells, so an RNA-derived ranking can be judged by a modality withheld from every retrieval method. We built two protein-based evaluators from those same measurements, changing only how they were summarized, and scored the identical two RNA-derived rankings against both (Fig.~\ref{fig:4}a). The first compared the mean protein responses of each query--candidate pair, and under it mean-signature retrieval outperformed energy-distance retrieval ($\SHAPEMEANORACLEMEAN$ versus $\SHAPEMEANORACLEENERGY$).

% FIGURE 4
\begin{figure}[htbp]
\centering
\includegraphics[width=\textwidth]{figures/fig4.pdf}
\caption{\textbf{External evaluation is not automatically neutral, and natural heterogeneity does not automatically imply decision reordering.}
Full definitions and robustness analyses are in Supplementary Notes~2 and~3.
\textbf{a}--\textbf{c}, One external modality judging the same rankings two ways.
\textbf{a}, One Perturb-CITE-seq experiment measures RNA and 20 surface proteins in the same 218,331 cells; the two RNA rankings and the control-centred protein responses are held identical while only the statistical form of the protein evaluator changes.
\textbf{b}, Spearman $\rho$ of three RNA-side scores against that evaluator, built once as a cosine between mean protein deltas and once as an energy distance between protein cell clouds; each bar is the median over $n=659$ queries of a per-query rank correlation across that query's candidate knockouts. Open bar, the response-magnitude control, which ranks candidates but compares no distributions.
\textbf{c}, How far energy exceeds the magnitude control under each form, pooled and within each immune condition (ratios 1.06 to 1.84); no uncertainty is shown, because the source records one correlation per condition per form.
\textbf{d}, The HIR-Bench boundary $\alpha^*=B/(A+B)$ separates settings in which the mean is sufficient from settings in which subpopulation structure can change the preferred candidate, without identifying them in biological data (Supplementary Note~2).
\textbf{e}, \textbf{f}, The same question in four glioblastoma tumours, each split by cell identity into disjoint malignant and myeloid compartments that share no cells and see the same drug and matched control.
\textbf{e}, Response-identity recoverability of within-patient drug-response structure; open bars, best unsupervised method; filled bars, supervised probe; median paired gap $+0.117$ (95\% patient-level cluster bootstrap $+0.115$ to $+0.135$, $n=36$ splits), against $+0.007$ in a constructed comparator.
\textbf{f}, Myeloid-response similarity against malignant-response similarity across those disjoint compartments, on which no retrieval was performed (Spearman $\rho=\PREMISEDISJOINT$, $n=18$ drug pairs, 15 of them from a single patient; dashed line, equal similarity).}
\label{fig:4}
\end{figure}

Against the full distributions of the same protein responses the ordering reversed: energy-distance retrieval reached $\SHAPEDISTORACLEENERGY$ against $\SHAPEDISTORACLEMEAN$ for mean-signature retrieval (Fig.~\ref{fig:4}b). Energy led under the distribution-based evaluator in all three immune conditions, and the change in which method wins occurred in two of the three. Changing only how the same held-out biological measurements were represented was therefore enough to reverse which retrieval strategy performed better.

Part of that reversal needs no distributional comparison. A scalar score that compares only the overall magnitude of the query and candidate responses rose from $\SHAPEMAGMEAN$ under the mean-based protein evaluator to $\SHAPEMAGDIST$ under the distribution-based one. Energy distance stayed ahead of that control under both forms, by $\SHAPEEDGEMEAN$ and $\SHAPEEDGEDIST$ and in every immune condition separately (Fig.~\ref{fig:4}c). Response magnitude therefore explains part of the reversal, while energy retains an additional association beyond that scalar channel. The same held-out measurements can favour different retrieval strategies depending on how their response structure is represented.

A controlled HIR-Bench construction confirms that subpopulation structure can change the preferred candidate where state-specific preferences conflict, although its analytical boundary did not transfer to the biological datasets (Fig.~\ref{fig:4}d; Supplementary Note~2).

We next asked whether state-dependent perturbation responses are recoverable in naturally heterogeneous tissue\cite{dataZhao}. In glioblastoma samples from ZhaoSims2021, the malignant and myeloid compartments of each tumour were analysed separately against compartment-matched controls. A supervised probe recovered within-patient drug-response identity at an accuracy of 0.923 against 0.777 for the best unsupervised procedure, a median paired gap of $+0.117$ against $+0.007$ in a constructed comparator (Fig.~\ref{fig:4}e; Supplementary Note~3).

Whether that structure is large enough to change which drug should be preferred is a separate question, and the tumour data constrain it. Across 18 within-patient drug pairs, myeloid-response similarity ranked malignant-response similarity at Spearman $\rho=\PREMISEDISJOINT$ across the two disjoint compartments, which share no cells (Fig.~\ref{fig:4}f). The pairs are not independent: 15 of the 18 come from one patient and the remaining three from one patient each, so no $P$ value is quoted and the coefficient describes this cohort rather than estimating a population value. Drug-response similarity was therefore largely preserved across distinct cellular compartments, leaving less room for state-specific structure to reorder candidates than the presence of heterogeneity alone would suggest. This measures preservation of drug ordering rather than retrieval performance, and it is the requirement the intervention-identification task below tests directly: cellular structure must not only exist and be recoverable, but must also change candidate preference.

\subsection*{Population information changes an intervention decision when candidate responses are observed}

A retrieval system that is to propose interventions must rank compounds whose effect in the query context has not been measured. We therefore ran the same comparison as a source-to-target intervention-identification task on Tahoe-100M\cite{Tahoe100M}: a held-out cell line's untreated cells and a desired treated population enter, and a ranking of a fixed \NCAND-compound library leaves, with the true intervention's rank as the outcome (\NQUERIES\ queries over \NCTX\ cell lines; Fig.~\ref{fig:5}a). Because the target is an observed response produced by one library compound, the task measures intervention identification rather than therapeutic utility. It was run twice, once with candidate responses observed within the query line and once with them predicted from the other \NFITCTX\ lines, which separates the value of population information from the cost of forecasting it.

% FIGURE 5
\begin{figure}[htbp]
\centering
\includegraphics[width=\textwidth]{figures/fig5.pdf}
\caption{\textbf{Population structure adds a small oracle retrieval gain beyond response magnitude, and most of it comes from a fifth of the queries.}
Source-to-target intervention identification on Tahoe-100M: \NQUERIES\ queries over \NCTX\ cell lines and a fixed \NCAND-compound library, with candidate responses \emph{observed} within the query line throughout and predicted from the other \NFITCTX\ lines in Fig.~\ref{fig:6}; \textbf{a} is schematic and energy distances use the unbiased U-statistic estimator.
\textbf{a}, A held-out line's untreated cells and a desired treated population enter and a ranking of all \NCAND\ compounds leaves; the target population is the query, not a candidate response.
\textbf{b}, Mean reciprocal rank for the three scoring rules of Fig.~\ref{fig:1}a, with 95\% bootstrap intervals: 0.945, 0.965 and 0.976, so the two steps are $+0.0201$ and $+0.0102$.
\textbf{c}, Top-1 decisions changed against each mean reference, ratios above the pairs; corrected, the population route puts the true drug first where the reference did not, and broken is the reverse.
\textbf{d}, The same $+0.0303$ pooled advantage split by whether the mean route left any room; on 3,110 of \NQUERIES\ queries it already ranks the true drug first at every seed.
\textbf{e}, Spearman correlation between that headroom and the population gain, against a null that reshuffles the population route's reciprocal ranks within each context, keeping both marginals, the $\mathrm{RR}\le1$ ceiling and the clustering.
\textbf{f}, Single-cell response detectability, separating a query drug's own treated cells from vehicle cells of the same line at up to 300 cells per arm; median accuracy 0.800 with the labels given and 0.555 without them, on a balanced two-class problem where 0.5 is chance.
\textbf{g}, Each interaction statistic against the two quantities it must not measure: the naive $1-\cos(r_1,r_2)$ runs at $-0.69$ with response magnitude and $-0.43$ with cells per arm, the cross-fitted $S_\mathrm{int}$ at $+0.66$ and $+0.10$.
\textbf{h}, Standardised coefficients from regressions of the population route's own reciprocal rank on each explanator, with the mean route's reciprocal rank as a covariate; the first three rows are one fit ($n=3{,}618$ queries in \NCTX\ cell-line clusters) and the naive cosine a separate one ($n=3{,}620$). Filled markers, coefficients whose interval excludes zero.}
\label{fig:5}
\end{figure}

With candidate responses observed, the three scoring rules of Fig.~\ref{fig:1}a again form a ladder. Mean reciprocal rank was 0.945 for direction-only mean cosine, 0.965 once response magnitude was restored and 0.976 for the energy distance between populations, so the step to the magnitude-aware mean is $+0.0201$ and the step from it to the population $+0.0102$ (Fig.~\ref{fig:5}b). That last step is not a property of population scoring as such. Of the two secondary population scores specified for this task, maximum mean discrepancy reproduces the energy result (MRR 0.975, $+0.0101$ over the magnitude-aware mean), whereas sliced Wasserstein distance reaches 0.922 and sits below both mean baselines ($-0.0230$ against direction-only matching, $95\%$ CI $-0.0301$ to $-0.0166$; $-0.0431$ against the magnitude-aware mean, $-0.0509$ to $-0.0359$), so two of the three population scores clear both mean baselines here and one does not. The advantage reaches individual decisions: against the magnitude-aware mean, population scoring put the true drug first on 2.0\% of queries where the mean route did not and broke a correct choice on 0.5\%, a fourfold asymmetry, and against direction-only matching the two rates are 5.6\% and 1.1\% (Fig.~\ref{fig:5}c).

The pooled oracle gain is limited by the headroom that mean matching leaves. On 3,110 of the \NQUERIES\ queries the direction-only mean route already ranks the true drug first at every seed, where the population route can only lose and does, by $-0.0036$ on average; the remaining 882 queries more than account for the pooled $+0.0303$, gaining $+0.1501$ (Fig.~\ref{fig:5}d). Headroom of that kind correlates with the population route's gain at Spearman $\rho=+0.79$, but a permutation null that preserves the ranking ceiling reproduces it at $+0.786$ ($+0.767$ to $+0.805$; Fig.~\ref{fig:5}e, Methods), so the correlation measures the room available for improvement rather than which queries benefit.

We next asked which measurable query properties explain the remaining advantage. A query drug's own treated cells are separable from vehicle cells of the same line at a median accuracy of 0.800 when the labels are given and 0.555 without them, on a balanced two-class problem where 0.5 is chance (Fig.~\ref{fig:5}f), and drug-by-state interaction was measured by a cross-fitted statistic that is zero in expectation under an additive null and largely free of the cell-count and magnitude dependence a naive response cosine carries (Fig.~\ref{fig:5}g; Methods). Conditional on the mean route's own reciprocal rank, response detectability is the only one of the three whose interval excludes zero ($+0.0091$, 95\% CI $+0.0060$ to $+0.0121$, $n=3{,}618$ queries in \NCTX\ cell-line clusters); neither the interaction share ($+0.0015$, $p=0.21$) nor its reproducibility across disjoint halves ($-0.0011$, $p=0.21$) is associated with the population route. Substituting the naive cosine for the cross-fitted statistic reverses the sign of the association ($-0.0096$, 95\% CI $-0.0167$ to $-0.0024$; Fig.~\ref{fig:5}h). Population-level information therefore has measurable value at the oracle, and little room in which to act.



% FIGURE 6
\begin{figure}[htbp]
\centering
\includegraphics[width=\textwidth]{figures/fig6.pdf}
\caption{\textbf{The oracle advantage does not survive forward prediction, and the deficit is a magnitude error rather than a distributional one.}
The same task, queries, library and split as Fig.~\ref{fig:5}, with candidate responses in the held-out line produced by a forward predictor fitted on the other \NFITCTX\ lines; panel \textbf{a} is schematic and intervals are 95\% bootstrap intervals over cell-line clusters.
\textbf{a}, The task with that one station replaced; every other element is drawn as in Fig.~\ref{fig:5}a.
\textbf{b}, Both routes travelling from observed to predicted candidates under the primary average-effect predictor; the mean route gives up 0.058 MRR and the population route 0.124, and the two cross.
\textbf{c}, $\Delta$MRR of the population route over direction-only mean matching, at the oracle and for six predictors; under an additive predictor the oracle gain inverts rather than attenuating.
\textbf{d}, The same predictions judged three ways: the population route's $-0.0356$ deficit against direction-only is matched by the magnitude-aware mean's $-0.0335$, and $-0.0020$ remains once magnitude is controlled.
\textbf{e}, Median drug-by-state interaction in each predictor's own output against the cross-fitted value measured in real cells (dashed); an additive predictor's is zero to float precision, an algebraic property of the model class.
\textbf{f}, A state-conditioned average-effect predictor is the first to generate real interaction without giving up mean-route accuracy (panel \textbf{e}), and it does not close the gap.
\textbf{g}, The same attrition counted over the 755 queries on which the oracle population route had any gain: 209 keep a gain under the primary average-effect predictor and 546 do not. Gains also appear where the oracle had none, 343 for that predictor and 1,165 for the Gaussian OT map, on a different denominator and, for that map, on a mean route reaching only 0.665.
\textbf{h}, The budget across both figures; the first two rows share the oracle MRR unit and are drawn to scale, and the last two are a rate and a signed gain and carry no width.}
\label{fig:6}
\end{figure}


\subsection*{The oracle advantage does not survive forward prediction}

Replacing the observed candidate responses with predicted ones removes the advantage and reverses its sign (Fig.~\ref{fig:6}a). The primary average-effect predictor costs the mean route 0.058 MRR and the population route 0.124, and the two cross (Fig.~\ref{fig:6}b). Under it the population route sits $-0.0356$ below direction-only mean matching (95\% CI $-0.0620$ to $-0.0115$), so the oracle gain does not attenuate but changes sign (Fig.~\ref{fig:6}c). None of the six forward predictors evaluated delivers a population advantage, and the Gaussian optimal-transport map is the only one whose point estimate is not negative ($+0.0027$, interval covering zero).

The dominant deficit is a magnitude error rather than a distribution-specific one. On the same predictions the magnitude-aware mean loses $-0.0335$ to direction-only matching, almost the whole of the population route's deficit, although it reads no distribution at all; against the magnitude-aware mean itself the population route sits only $-0.0020$ below ($-0.0036$ to $-0.0005$; Fig.~\ref{fig:6}d). An additive predictor translates the query line's own vehicle cells by an effect averaged over the other \NFITCTX\ lines, so it reproduces that population's heterogeneity exactly while generating no drug-by-state interaction at all, an algebraic property of the model class rather than a training failure (Fig.~\ref{fig:6}e; Supplementary Note~4).

A state-conditioned average-effect predictor is the first to generate real interaction while preserving the mean route's MRR (0.8883 against 0.8876), and it still leaves the population route behind ($-0.0314$, $-0.0584$ to $-0.0076$; Fig.~\ref{fig:6}f). Generating state-specific interaction is therefore necessary for the population route but not sufficient. The same attrition is visible query by query: of the 755 queries on which the oracle population route had any gain, the average-effect predictor keeps one on 209 and loses it on 546 (Fig.~\ref{fig:6}g).

Together the two arms locate where population-level information enters and where it is lost (Fig.~\ref{fig:6}h). With responses observed it is real, it reaches a ranking decision, and it is bounded by how good mean signatures already are. With responses predicted it does not survive: the predictors that best preserved mean-level retrieval were additive and generated no state-specific perturbation response, and introducing state conditioning restored that interaction without restoring a population-level advantage.

\section*{Discussion}

The results trace one chain of attrition rather than a checklist. Response populations carry retrievable information; a smaller part of it exceeds response magnitude; a smaller part again changes which candidate is preferred; and almost none of that survives once candidate responses must be predicted rather than observed. The binding constraint changes along that chain: at the oracle it is the ranking headroom that mean matching leaves, because mean matching already ranks the true compound first on most queries, and once candidate outcomes are unobserved it is the fidelity of the forward predictor.

Single-cell perturbation atlases\cite{ref3,dataSciPlex3,Tahoe100M} change the representation available for drug retrieval: a perturbation can now be described by a population of cellular responses rather than by a single averaged signature. Across the datasets and evaluators examined here, retaining the full population increased the information available to the retriever, but only a smaller component beyond response magnitude reached the ranking, and it did so under response matching, where Fig.~\ref{fig:2}e shows that it adds nothing to the magnitude-aware mean under the regret criterion; under mechanism-of-action recovery the whole population advantage does not survive (Fig.~\ref{fig:3}a,b,e). The central distinction is therefore not mean versus distributional retrieval, but information that is available versus information that changes the decision. Classical perturbation signatures and single-cell populations are points on a continuum of retained response structure, and the relevant question is how much of that structure the decision at hand requires.

This distinction has direct implications for artificial intelligence virtual cells\cite{AIVirtualCell,VirtualCellChallenge}. A virtual cell that accurately predicts an average perturbation response can still fail to provide the information needed for population-level intervention selection. A model that applies the same perturbation effect to every cell can preserve heterogeneity in the starting population yet, by construction, cannot generate state-specific differential responses; the limit is exact for strictly additive predictors, although more flexible latent models are not algebraically constrained to it (Supplementary Note~4). A downstream population-aware retriever cannot recover variation that the predictor never generated. Reproducing state-dependent differential response is therefore necessary to exploit that particular source of population information, and it is not sufficient: introducing state conditioning restored measurable drug-by-state interaction while preserving mean-level performance, yet did not restore a population-level retrieval advantage (Fig.~\ref{fig:6}e,f). The same experiment identifies where to look next, because the magnitude-aware mean loses almost exactly what the population route loses, so the dominant remaining error is predicted response magnitude rather than population shape (Fig.~\ref{fig:6}d). That is a concrete target for virtual-cell modelling.

Two limits bound these conclusions. The response and interaction diagnostics examined here did not yield a reliable query-level rule for when population scoring helps, and state-dependent differential response is only one of the forms population structure could take, so a gain arising from another form is not excluded. These results also suggest a practical principle for evaluating population-aware retrieval. The statistical representation of an external readout can itself change which method ranks higher\cite{ref17,MetricGaming}, so validation should ask whether the added cellular information changes an outcome relevant to the intended use, and should compare against simpler channels such as response magnitude\cite{ref13} before attributing a gain to population structure itself. The functional readouts analysed here remain indirect and do not establish therapeutic benefit. More broadly, a richer cellular measurement justifies a richer retrieval method only when the information it adds survives to the decision that method exists to make.


%% =====================================================================
%%  METHODS
%% =====================================================================
\section*{Methods}

\subsection*{Study design and data resources}

This study asks when retaining the full population of single-cell perturbation responses improves retrieval beyond a conventional mean expression signature. The analyses were therefore organized around four distinct roles: observed-population retrieval, evaluation against biological and functional criteria, measurement of naturally occurring response heterogeneity, and large-scale assessment of the conditions under which population-level information can be useful. These roles were analysed separately rather than pooled into a single benchmark.

The primary observed-population retrieval analyses used SciPlex3\cite{dataSciPlex3}, comprising 276,325 cells across three cancer cell lines (A549, K562 and MCF7), 188 compounds shared across the three contexts and a 2,000-gene expression space. SciPlex3 supports both within-context retrieval and controlled cross-context constructions because the same compound library is measured across all three cell lines. Frangieh melanoma Perturb-CITE-seq\cite{dataFrangieh} spans Control, IFN$\gamma$ and Co-culture conditions with 239 CRISPR knockouts. Both analyses read the same RNA matrix under two filters: the retrieval analysis uses the 218,023 cells left after dropping nine knockout targets with fewer than 50 cells, and the evaluator-shape experiment uses all \SHAPECELLS\ cells of the RNA--protein barcode intersection, of which 215,915 contribute to at least one condition mean once the 60-cell floor per knockout is applied. We use this dataset as an independent perturbation-retrieval setting and, because matched surface-protein measurements are available, for the evaluator-shape experiment described below, in which that protein modality is held out from every retrieval method and used only to judge. It is a genetic-perturbation test of the same retrieval formulation rather than a drug-retrieval benchmark. CD34$+$ haematopoietic progenitor data (GSE306429)\cite{dataCD34} contain 33,984 cells and 36 compounds and are used as a negative-control setting for population-level retrieval.

Three additional resources serve roles outside the primary retrieval benchmark. ZhaoSims2021\cite{dataZhao} contains drug-treated patient-derived glioblastoma tissue and is used to test whether state-dependent responses and their recoverability are present in naturally heterogeneous tissue; no retrieval performance is reported from this cohort. Tahoe-100M\cite{Tahoe100M} serves two distinct roles: it provides larger-scale measurements of the same information requirements under intrinsic within-cell-line heterogeneity, and it is the substrate for the source-to-target intervention-identification benchmark of Figs.~\ref{fig:5} and~\ref{fig:6}. We analyse plate 3 only, which comprises 4,158,278 cells across 50 cell lines and 93 compounds, using the released 2,304-gene log1p-transformed representation without refitting the preprocessing pipeline. Finally, GDSC2 dose--response profiles provide an external functional readout for drug--drug similarity and are joined to the SciPlex3 compound panel independently of the transcriptomic retrieval scores.

Observed responses, predicted responses, external functional measurements and natural-tissue analyses therefore answer different questions in this study. Results are not combined across these resources as though they represented independent replicates of one benchmark. Dataset scale, inclusion criteria, experimental role and query construction are summarized in Supplementary Table~2.


\subsection*{Preprocessing and construction of perturbation-response populations}

Single-cell expression analyses were performed in log-transformed expression space using the processed gene sets defined for each dataset. SciPlex3, Frangieh and CD34$+$ were analysed in 2,000-gene representations. For SciPlex3 and
Frangieh these genes were selected in this study rather than inherited from the released objects:
counts were normalized to a total of $10^4$ per cell and log1p-transformed, 2,000 highly variable
genes were chosen by the Seurat criterion\cite{Seurat} on a subsample of cells, and the final matrices were then
rebuilt from raw counts on those genes using all-gene size factors so that normalization is applied
once. SciPlex3 was additionally downsampled to at most 120 cells per compound, cell line and dose,
and 3,000 control cells per cell line. The CD34$+$ tensor was received in processed form and the feature-selection
procedure applied to it is not documented by the source. For ZhaoSims2021, counts were normalized to a total of $10^4$ per cell, log1p-transformed and reduced to 2,000 highly variable genes. Tahoe-100M was used in the released preprocessed form (2,304 highly variable genes, log1p transformed), with no additional feature selection or representation fitting. Dataset-specific filtering and minimum-cell requirements for individual analyses are given in the corresponding Methods sections below.

For a perturbation $d$ measured in cellular context $c$, let
\[
X_{d,c}=\{x^{(d,c)}_1,\ldots,x^{(d,c)}_{n_{d,c}}\}
\]
denote its treated single-cell expression profiles and let
\[
\mu_{0,c}
=
\frac{1}{n_{0,c}}
\sum_{j=1}^{n_{0,c}}x^{(0,c)}_j
\]
denote the mean of the matched control population. The conventional perturbation signature is the control-referenced mean response
\[
\bar d_{d,c}
=
\frac{1}{n_{d,c}}
\sum_{i=1}^{n_{d,c}}x^{(d,c)}_i
-
\mu_{0,c}.
\]
Mean-signature retrieval operates only on $\bar d_{d,c}$.

Population-level retrieval retains the corresponding control-referenced cellular responses,
\[
P_{d,c}
=
\left\{
x^{(d,c)}_i-\mu_{0,c}
\right\}_{i=1}^{n_{d,c}},
\]
so that the mean of $P_{d,c}$ is exactly the conventional mean perturbation signature while variation among its cells remains available to the retrieval score. Both arms are therefore built from the same treated cells, the same matched control, and the same query and candidate libraries. The energy distance and the mean $L_2$ depend only on distances between cells and are unchanged by subtracting a control that a query and its candidates share, so for those two scores the referencing is a change of origin rather than of comparison; the mean cosine is not invariant in this way and is computed on the control-referenced means. Where a query and its candidates instead carry different matched controls, as the two single-mode candidates of the two-context mixture task do, that difference enters the mean scores and not the population distance.

Controls were matched to the biological context in which the perturbation was measured. For analyses that explicitly compare responses between cellular states or compartments, each state is instead referenced to its own matched control population so that pre-existing differences between states cannot be mistaken for differential perturbation responses. This state-specific construction is used for the differential-response analyses in predicted populations, patient-derived glioblastoma and Tahoe-100M and is defined in the corresponding sections below.

Cell sampling was performed only after the response populations had been defined. Minimum cell numbers, per-condition caps and random seeds differ by experiment and are reported with the corresponding analysis rather than imposed globally. Protein measurements used for the evaluator-shape experiment were processed separately, including barcode-based RNA--protein matching and protein-specific normalization, as described below.


\subsection*{Population-to-population retrieval and signature baselines}

For a query response population $Q$ and a library of candidate response populations
$\{P_d\}_{d=1}^{D}$, retrieval assigns each candidate a similarity or distance to the
query and orders the candidate library accordingly. Population-level methods operate
directly on the empirical response populations defined above, whereas signature-based
methods first reduce each population to its control-referenced mean response. Query
construction and candidate libraries are kept identical when comparing the two
representations.

\paragraph{Population-level retrieval.}
We evaluated energy distance\cite{EnergyDistance} (the primary population score), maximum
mean discrepancy\cite{MMD}, sliced Wasserstein distance\cite{SlicedWasserstein}, and the
mean and worst-case limits of a subpopulation-coverage score. All rank candidates from
smallest to largest distance. For a candidate population $P=\{x_1,\dots,x_m\}$ and a query
population $Q=\{y_1,\dots,y_n\}$ the energy distance is estimated as
\[
\widehat{\mathcal{E}}(P,Q)
=
\frac{2}{mn}\sum_{i,j}\lVert x_i-y_j\rVert
-
\frac{1}{m(m-1)}\sum_{i\neq i'}\lVert x_i-x_{i'}\rVert
-
\frac{1}{n(n-1)}\sum_{j\neq j'}\lVert y_j-y_{j'}\rVert ,
\]
the Euclidean norm itself rather than its square. Excluding the self-pairs makes this the
unbiased U-statistic form; the biased V-statistic divides the two within-population terms by
$m^2$ and $n^2$ instead, which adds each population's own mean within-population distance
divided by its cell count and therefore shifts candidates of different size or dispersion by
different amounts. Every energy distance reported here is the U form except the PCA-latent
population score of Fig.~\ref{fig:2}a, which is computed by a separate routine in the V form;
the two forms are compared arm by arm in the code release. The population distances are computed
on normalized expression rather than on control-referenced responses, whereas the mean-signature
scores subtract the matched control mean. Fixed cell caps, penalties for undersized populations
and the implementation settings for secondary scores are reported in Supplementary Methods.

\paragraph{Subpopulation-coverage score.}
The query is partitioned into $K$ components ($k=2$ throughout) and each candidate cell is
assigned to the nearest query component. In the observed-population experiments the partition is
the construction label of the two response modes and assignment is by Euclidean nearest centroid;
in the partial-observation experiments, where no such label exists, the partition is $k$-means on
the query cells with three initializations and a fixed seed. Writing
$d_k$ for the base distance between the query cells and the candidate cells of
component $k$, the coverage distance and its ranking score are
\[
D_\beta=\frac{1}{\beta}\Big(\log\sum_{k=1}^{K}e^{\beta d_k}-\log K\Big),
\qquad
S_\beta=-D_\beta .
\]
The base distance $d_k$ is one of energy, MMD or sliced Wasserstein, so coverage is a
family built on the same kernels rather than an independent metric. A candidate that
populates no cell of a query component receives the penalty distance for that
component, which is how the score charges a candidate for leaving part of the query
uncovered.

As $\beta\rightarrow0$, $D_\beta$ averages subpopulation distances; as
$\beta\rightarrow\infty$, it approaches the worst-matched component: $0.7125$ and
$1.300$ for the four per-subpopulation distances drawn in Fig.~\ref{fig:1}h. For $K=1$
it equals the global distance. Coverage-mean and coverage-worst are these two limits; numerical
implementation details are given in Supplementary Methods.

\paragraph{Mean-signature baselines.}
For each response population $P_{d,c}$, the corresponding mean signature is
$\bar d_{d,c}=\operatorname{mean}(P_{d,c})$. Mean-signature retrieval ranks candidates
using either cosine similarity or Euclidean distance between these vectors. The
implemented CMap-style cosine baseline applies cosine similarity to the same
control-referenced mean differential-expression vectors and is therefore numerically
identical to mean-cosine retrieval (Spearman $\rho=1.000$ over 54,180 query--candidate
scores). The canonical L1000 weighted connectivity score
(WTCS)\cite{ref2} is included as a separate signature-based baseline. WTCS uses rank-based
enrichment of the mean-collapsed expression response and is related to, but not
algebraically identical to, the cosine implementation used here; it gives a different
Hit@1 (\HITWTCS).

To test whether the comparison depended on the original gene space, we also evaluated
a PCA representation. One basis was fitted per retrieval query, on the query cells pooled with
the cells of all 43 candidate populations of that query, the held-out correct candidate included,
uniformly subsampled without replacement to 4,000 cells at a fixed seed and reduced to 30
components. The fit is therefore transductive but unsupervised: no candidate identity, ground-truth
label or subpopulation label enters it, the correct candidate contributes on the same terms as the
other 42, and the query cells are disjoint from the correct candidate's cells by construction.
Control cells do not enter the fit, each control having already been collapsed to a mean vector
that is projected through the fitted basis. The same comparison was then performed either between
latent-space means or between the corresponding latent response populations.

\paragraph{Connection between mean and population retrieval.}
To formalize the relationship between the two representations, we write a response
population as
\[
X=\mu+\lambda\epsilon,
\qquad
E[\epsilon]=0,
\]
where $\mu$ is the mean response, $\epsilon$ represents within-population variation and
$\lambda$ controls its magnitude. As $\lambda\rightarrow0$, the response population
collapses to a point mass at $\mu$, and population comparison therefore loses all
information beyond the mean. The mean \emph{representation} is the endpoint of this
variance-sensitivity construction. The score obtained at that endpoint is a distance
between two mean vectors: for the energy distance between two point masses it is
exactly $2\lVert\mu_Q-\mu_d\rVert_2$, which is the magnitude-aware mean score and not
the cosine. Which score reads the mean is therefore a separate choice from how much of
the population is retained, and the cosine is a further coarsening of the same
representation that discards $\lVert\mu_Q-\mu_d\rVert_2$ altogether. Figure~\ref{fig:2}a
measures what that coarsening costs.

The continuum refers to the amount of within-population response structure available to the
scorer, not to an equivalence between otherwise different scoring functions: WTCS remains a
distinct rank-based signature score, and energy, MMD and sliced Wasserstein retain their own
distance definitions.


\subsection*{Response-matching and partial-observation evaluation}

\paragraph{Observed-population retrieval.}
When complete candidate response populations were available, retrieval was evaluated over
heterogeneous queries. Each query is a population built by combining two response modes at mixing
weight $\alpha\in\{0.5,0.7,0.9\}$, and the correct answer is the candidate that reproduces both
modes at the same weight, built from the held-out half of the same cells so that no cell is shared
between a query and its answer. Two further candidates carry one mode each, and 40 distractor
populations complete a 43-candidate library.

The three tasks differ in what the two modes are. In the controlled SciPlex3 task they are the
HDAC- and JAK-associated response sources within a single cell line at 10\,$\mu$M, each pooling
every compound annotated to that class, so the two single-mode candidates are HDAC-only and
JAK-only populations and each distractor is an individual compound of another class. In the
cross-line and Frangieh tasks they are one perturbation observed in two biological contexts, two
cell lines or Control and IFN$\gamma$, so the two single-mode candidates are that perturbation
observed in one context only and each distractor is a different perturbation mixed across the same
two contexts at the same weight; cross-line perturbations pool all measured doses. The two
single-mode candidates are each referenced to their own single-context control, whereas the correct
answer and every distractor are referenced to the $\alpha$-mixed control.

Query perturbations in these two tasks are not arbitrary. Among perturbations that pass a
split-half reliability floor in both contexts and a cross-context response cosine below 0.9, the
analysis uses the most divergent: 12 drugs for each cell-line pair and 6 Frangieh knockouts. These
two tasks therefore measure retrieval where cross-context divergence is largest rather than on a
representative perturbation panel.

Populations contain 400 cells in the controlled task and 200 in the other two, sampled without
replacement where the source pool allows it and with replacement otherwise. Within one task
setting every query is drawn from a single fixed split of that setting's cells and differs only in
the perturbation or mechanism sampled, in $\alpha$ and in the resampling seed, so queries within a
setting are not independent replicates. The controlled and cross-line tasks use ten resampling
seeds and the Frangieh task five.

Hit@1 is reported as an unweighted macro-mean over the seven task--setting combinations, so a
30-query combination counts as much as a 360-query one. Weighting instead by query count gives
\HITENERGYW\ for energy distance and \HITMEANW\ for mean cosine, and PCA-space retrieval reaches
\HITPCADIST\ against \HITPCAMEAN\ for the corresponding mean-based score.

Hit@1 is the fraction of queries in which no candidate scores strictly higher than the both-mode
candidate; its rank is one plus the number of strictly higher-scoring candidates, and the same
definition underlies Hit@5, MRR and nDCG@10. That convention would resolve an exact tie in the
both-mode candidate's favour, but no tie arose: across the 1,260 queries and all nine scorers the
both-mode candidate is never tied with another candidate at any rank, so every Hit@1 reported here
is identical under optimistic, pessimistic and fractional tie handling. The partial-observation protocol below is scored differently: it ranks by a label-blind stable
sort and its Hit@1 is the best rank attained by any member of the query's mechanism or target set,
not the rank of a single correct candidate. Ranking a single-mode candidate first counts as a miss, including in the cross-line and
Frangieh tasks, where that candidate carries the correct perturbation observed in one of the two
contexts. The criterion therefore scores the composition of the retrieved population as well as the
mechanism or perturbation it was drawn from.
Comparisons between population-level and mean-signature retrieval use the same queries and
candidate libraries. Aggregate Hit@1 values are reported as unweighted macro-means over the seven
task--setting combinations used in the observed-population analysis; query-weighted averages are
reported as a sensitivity analysis so that the largest task does not determine the primary summary.

\paragraph{Partial-observation protocol.}
We additionally evaluated retrieval when only part of each candidate response population was
available to the retriever. This setting was designed to separate the information used to construct
a ranking from the criterion subsequently used to evaluate that ranking. For each query, the
population-level and mean-signature methods rank the same candidate library once; the resulting
rankings are then retained unchanged while different evaluation criteria are applied.

Three splits construct the partial observation, all on SciPlex3, and in each the hidden material
never enters the candidate library\cite{Leakage}. In \emph{leave-drug-out},
one drug is hidden, the query is that drug's own treated population (up to 200 cells) and the
library is every other drug in the context. In \emph{leave-MoA-out}, an entire annotated
mechanism class is removed from the library and the query is one drug drawn from the removed
class, so recovery requires a same-target surrogate among the remaining drugs. In
\emph{partial-library}, a fraction of the drugs is treated as measured, the query is built from
the pooled responses of the unmeasured drugs, and the retriever selects the next batch of
candidates. In all three, query states are obtained by $k$-means with $k=2$ on the query cells
and the minority state is the smaller cluster; candidate populations are capped at 120--150
cells per drug.

The response-matching criterion is an energy-based welfare regret. Each candidate's welfare is
the aggregate over query states $s$ of its negative energy distance to that state,
$U(d)=\operatorname{agg}_s\big(-\mathcal{E}(P_d,Q_s)\big)$ with the aggregator taken as the
minimum in the worst-case form, and the regret of the selected candidate $\hat d$ is
\[
\mathrm{Regret}(\hat d)=\max_{d\in\mathcal{L}}U(d)-U(\hat d)\;\ge\;0,
\]
where $\mathcal{L}$ is the observed library, so the reference is the best candidate actually
available rather than an unobservable optimum. Minority-state coverage is reported on the same
rankings as the cosine between the candidate's mean response and the mean response of the
query's minority state, rescaled from $[-1,1]$ to $[0,1]$. Regret is zero when the top-ranked
candidate is optimal under that objective and increases as the selected candidate becomes less
favourable. We report regret reduction relative to mean-signature retrieval, such that a positive
value indicates lower regret under this response-matching criterion. This evaluation is intentionally
closely aligned with the population-level information used by the retriever and is therefore
interpreted as evidence of response fidelity rather than as an independent functional outcome.

\paragraph{Task-proximal biological evaluation.}
The same fixed rankings were next evaluated using mechanism-of-action (MoA) recovery and
minority-state coverage. MoA recovery was quantified by normalized discounted cumulative gain
(nDCG)\cite{nDCG}, which rewards candidates sharing the query perturbation's annotated mechanism
when they occur near the top of the ranked list. Mechanism recovery is scored on the \NMOA\ leave-drug-out queries, that is on the
\NALL\ partial-observation queries less the 120 leave-MoA-out and 45 partial-library queries,
whose candidate libraries are built without a member of the query mechanism class and for which
recovery is therefore undefined. Selection is by split type. Within the \NMOA, 21 queries happen
to hold no mechanism-class member either and score exactly zero under every method, so they enter
the comparison as ties rather than as a difference; dropping them leaves the median difference at
0.000 and moves the mean from $\CLASSBMEAN$ to $-0.025$ and the decreased and increased shares
from \CLASSBWORSE\ and \CLASSBBETTER\ to 42.3\% and 37.0\%.

Minority-state coverage asks whether the top-ranked candidate better preserves the response of the
less represented cellular state in the controlled heterogeneous-population construction. As with
MoA recovery, this criterion is applied to rankings that have already been generated; the evaluator
does not alter the query, the retrieval score or the candidate ordering. Differences between
population-level and mean-signature retrieval are therefore attributable to how the same rankings
score under the biological criterion rather than to refitting or re-ranking.

\paragraph{Cross-setting projection.}
To examine whether the minority-state result extends beyond the controlled partial-observation
setting, we evaluated the top-ranked population-level and mean-signature candidates across 239
retrieval task instances spanning natural and constructed settings. These comprised 90 constructed
cross-cell-line SciPlex3 tasks, 90 within-cell-line SciPlex3 tasks, 23 Frangieh tasks, 12
CD34$+$ tasks and 24 tasks using predicted candidate populations. For each task we report the difference in minority-state coverage between the
population-level and mean-signature selections.

The 239 tasks are treated as a distribution of task-level differences rather than as independent
biological replicates of a single experiment. In particular, the constructed SciPlex3 tasks share
data source, preprocessing and candidate structure, and the three SciPlex3 cell lines constitute
leave-one-cell-line-out variation within one experiment rather than leave-one-dataset-out
generalization. Thresholded counts of ``distributionally dominant'' tasks are therefore not used as
a primary result; the continuous difference in coverage is reported instead.


\subsection*{External functional evaluation using GDSC2}

To test whether transcriptomic retrieval rankings capture information that extends beyond the
single-cell expression measurements used to construct them, we evaluated the SciPlex3 rankings
against drug-response profiles from the Genomics of Drug Sensitivity in Cancer\cite{GDSC,GDSCLandscape},
using the GDSC2 screening set (release 8.5; Data availability).
The external criterion was defined as similarity between the functional response profiles of two
drugs rather than as absolute drug potency, because the retrieval task asks which candidate produces
a response most similar to the query rather than which candidate is most cytotoxic.

SciPlex3 compounds were matched to GDSC2 by normalized drug name, with substring matching used
only as a fallback. After removal of one duplicate mapping and restriction to the GDSC2 screening
platform, 35 compounds remained. GDSC1 and GDSC2 measurements were not pooled because they derive
from different assay and dose-range settings. For each matched compound, its functional profile was
the vector of fitted dose--response AUC values across GDSC2 cell lines. The three cell lines used for
the transcriptomic retrieval experiments, A549, K562 and MCF7, were excluded from these profiles,
leaving \CTXLINES\ cell lines,
so the functional evaluation was performed in biological contexts not used to construct the RNA
rankings.

Cell lines differ systematically in their overall sensitivity to drug treatment, which can induce
positive correlations between otherwise unrelated compound profiles. We therefore removed this
shared cell-line effect before calculating drug--drug similarity. For cell line $l$ and compound
$d$, the centred response was
\[
\widetilde{\mathrm{AUC}}_{d,l}
=
\mathrm{AUC}_{d,l}
-
\frac{1}{|\mathcal D_l|}
\sum_{d'\in\mathcal D_l}\mathrm{AUC}_{d',l},
\]
where $\mathcal D_l$ contains all GDSC2 compounds measured in cell line $l$, rather than only the
35 compounds overlapping SciPlex3. Functional similarity between query drug $q$ and candidate
drug $c$ was then defined as
\[
F(q,c)
=
\rho_{\mathrm{S}}
\left(
\widetilde{\mathrm{AUC}}_{q,\bullet},
\widetilde{\mathrm{AUC}}_{c,\bullet}
\right),
\]
where $\rho_{\mathrm{S}}$ denotes Spearman correlation across cell lines shared by the two
compounds. Drug pairs measured in fewer than 100 common cell lines were excluded. This yielded
591 usable drug pairs, with a median of 909 shared cell lines per pair. The corresponding
uncentred functional similarity was retained as a sensitivity analysis.

For each SciPlex3 cell line, retrieval was evaluated in a leave-one-drug-out design. The
10\,$\mu$M response population of the held-out compound served as the query, and the remaining
matched compounds formed the candidate library, giving 34--35 candidates per query. Each RNA-based
retrieval rule produced a candidate ranking without access to GDSC2 data. Agreement between that
ranking and the external functional ordering $F(q,c)$ was quantified by Spearman correlation
across candidates.

The primary RNA-derived rankings were energy-distance population retrieval and control-referenced
mean-cosine retrieval. To test whether apparent functional agreement required cell-to-cell
distributional information, we additionally used a response-magnitude matching score,
\[
M(q,c)
=
-\left|
\|\bar d_q\|-\|\bar d_c\|
\right|,
\]
which prefers candidates whose overall transcriptional response magnitude is similar to that of
the query but uses no within-population structure. A query-independent magnitude score
$\|\bar d_c\|$ was retained as a diagnostic control. An uncorrected mean-cosine score based on
treated means without matched-control subtraction and a potency-matching diagnostic were analysed
as additional sensitivity checks rather than as primary retrieval baselines.

To estimate the component of the energy-based association not explained by response magnitude, we
computed a partial rank correlation with the magnitude-matching channel removed. Retrieval scores
and functional similarities were rank-transformed, each was residualized with respect to
$M(q,c)$, and the correlation between the two residuals was reported as the magnitude-adjusted
association.

Per-query Spearman correlations are summarized by their median within each SciPlex3 cell line.
Queries from the same cell line share the same candidate library and are therefore not treated as
independent biological replicates. The three cell lines provide the higher-level analysis unit;
query-level paired tests are reported only as descriptive analyses where indicated.


The comparison between energy retrieval and the response-magnitude scalar is pre-specified as a
pooled test across the three SciPlex3 cell lines; per-line ratios are given in Supplementary
Methods.

\subsection*{Protein-based evaluator-shape experiment}

Frangieh Perturb-CITE-seq\cite{dataFrangieh} provides matched RNA and surface-protein measurements from the same
cells, allowing the statistical form of an external evaluator to be changed while holding its
biological measurement fixed. RNA and protein matrices were joined by cell barcode before
analysis; row position was not used because the released modality-specific matrices differ in
cell number and ordering. Of the 24 measured antibodies, four isotype controls were removed after normalization,
leaving 20 surface-protein features. Antibody-derived tag counts\cite{CITEseq} were centred-log-ratio normalized within each cell,
$y_{ij}=\ln(1+x_{ij})-m^{-1}\sum_{k}\ln(1+x_{ik})$ for cell $i$ and antibody $j$, with a
pseudocount of one, a natural logarithm and the per-cell mean taken over all $m=24$ measured
antibodies including zero counts and the four isotype controls; normalization therefore precedes
isotype removal. This is the centred log-ratio of the $(1+x)$ vector and differs from
implementations that exclude zero entries from the geometric mean.

All analyses were performed separately within the Control, IFN$\gamma$ and Co-culture conditions
to avoid mixing immune context with perturbation identity. Knockouts represented by at least
60 barcode-matched cells were retained, and populations were subsampled to at most 250 cells
using a fixed random seed.

\paragraph{Protein-based evaluators.}
For knockout $k$ in immune condition $c$, let
$\bar d^{\,\mathrm{prot}}_{k,c}$ denote its mean control-referenced protein response and let
$D^{\mathrm{prot}}_{k,c}$ denote the corresponding set of per-cell protein responses. The
mean-based evaluator was defined as

\[
E_{\mathrm{mean}}(q,d)
=
\cos
\left(
\bar d^{\,\mathrm{prot}}_{q,c},
\bar d^{\,\mathrm{prot}}_{d,c}
\right).
\]
The population-based evaluator was defined from the same cells and protein measurements as

\[
E_{\mathrm{pop}}(q,d)
=
-
\mathcal{E}
\left(
D^{\mathrm{prot}}_{q,c},
D^{\mathrm{prot}}_{d,c}
\right),
\]
where $\mathcal{E}$ is the energy distance. The negative sign converts distance to similarity
so that larger values indicate a closer match under both evaluators.

For both definitions, protein responses were referenced to the matched control population from
the same immune condition. Thus,

\[
\bar d^{\,\mathrm{prot}}_{k,c}
=
\operatorname{mean}
\left(
X^{\mathrm{prot}}_{k,c}
\right)
-
\mu^{\mathrm{prot}}_{0,c},
\]
and

\[
D^{\mathrm{prot}}_{k,c}
=
\left\{
x^{\mathrm{prot}}_{i,k,c}
-
\mu^{\mathrm{prot}}_{0,c}
\right\}_{i=1}^{n_{k,c}}.
\]
Using the same matched control for the two evaluator forms ensures that their difference arises
from whether protein responses are summarized by their mean or retained as a cellular population,
rather than from a change in biological reference.

\paragraph{Fixed RNA-derived rankings.}
The RNA-derived rankings were computed once and then evaluated against both protein criteria.
The primary rankings were energy-distance population retrieval and control-referenced
mean-cosine retrieval. The query populations, candidate libraries and RNA-derived candidate
scores were identical under the two protein evaluators, and no protein measurement was available
to any retrieval method. The experiment therefore changes only the statistical representation
used to judge the same RNA-derived ranking.

\paragraph{Response-magnitude control.}
Because a population distance can also encode differences in overall response magnitude, we
included a query-dependent scalar control that uses no cell-to-cell distributional information.
For query $q$ and candidate $d$, magnitude matching was defined as

\[
M(q,d)
=
-
\left|
\left\|
\bar d^{\,\mathrm{RNA}}_{q,c}
\right\|
-
\left\|
\bar d^{\,\mathrm{RNA}}_{d,c}
\right\|
\right|.
\]
This ranking favours candidates with a transcriptional response magnitude similar to that of the
query while discarding all within-population structure. The same magnitude ranking was evaluated
against both protein criteria.

For each query, agreement between an RNA-derived ranking and a protein-based evaluator was
quantified by Spearman correlation across that query's candidate knockouts. The statistic is
therefore a rank correlation over candidates, computed once per query, and never a correlation over
pooled query--candidate pairs. The headline value for each ranking and each evaluator form is the
median of those per-query correlations over all 659 queries. The same medians were also taken
within each of the three immune conditions, which are reported separately and treated as the
higher-level analysis unit; queries within a condition share candidate structure and were not
interpreted as independent biological replicates.


% Detailed SciPlex3 predict-then-rank and response-identity diagnostics are retained in
% Supplementary Notes 1 and~4 and Supplementary Table 4; no main-text claim depends on them.
\iffalse
\subsection*{Differential-response analysis of predicted populations}

Predictor diagnostics were used only to ask whether candidate generation preserves
state-specific response variation, not to rank perturbation-prediction architectures.

\paragraph{Predict-then-rank response sources.}
Predicted candidate populations were generated by applying modelled effects to real control cells
in the query context. This SciPlex3 diagnostic is an in-sample ceiling on candidate generation;
predictor families and full specifications are provided in Supplementary Table~4.

\paragraph{Within-context definition of differential response.}
Differential response was measured within a single cellular context so that variation between
contexts could not be mistaken for state-specific perturbation effects. For each SciPlex3 context,
up to 400 untreated control cells were represented by 30-component PCA and partitioned into two
cellular states using $k$-means with $k=2$. Subpopulation sizes and silhouette scores were retained
as diagnostics of the underlying state partition.

For perturbation $d$, each predictor was applied separately to the control cells belonging to the
two states. The induced response of state $k$ was defined as

\[
r_{d,k}
=
\operatorname{mean}
\left(
\widehat{X}_{d,k}
\right)
-
\operatorname{mean}
\left(
X_{0,k}
\right),
\]
where $X_{0,k}$ denotes the untreated cells assigned to state $k$ and
$\widehat{X}_{d,k}$ denotes their predicted post-perturbation profiles. Differential response was
summarized by the cosine between the two state-specific response vectors,

\[
R_d
=
\cos
\left(
r_{d,0},
r_{d,1}
\right).
\]
A value near 1 indicates that the two cellular states are predicted to respond in nearly the same
direction, whereas lower values indicate increasingly state-specific perturbation responses. Real
treated cells from the same context, partitioned according to the same control-derived state
definition, were analysed identically and provide the observed reference. $R_d$ is an alignment
between two response directions and is used here to characterize predictor output, where the
additive limit pins it to exactly 1. Drug-by-state interaction in the retrieval analyses is measured
by the cross-fitted statistic defined below, which does not share the small-sample behaviour of a
cosine between two noisy means.

The structure diagnostics reported for the constructed cross-line mixtures are ratios to the
real population each predicted population replaces. Subpopulation variance is 0.138 in the real
populations and 0.142 under both the average-effect predictor and nearest-neighbour transfer,
against 0.462 under the linear-latent model. Taken as ratios to the real value, the average-effect
and nearest-neighbour predictors reproduce subpopulation variance, response diversity and isotropy
to within 0.99 to 1.04. The induced response cosine between the two cellular states is 0.014 in
the real populations against 0.186, 0.239 and 0.366 for the three predictors.

\paragraph{Additive-response limit.}
The limiting behaviour of an additive predictor follows directly from its model form. If the same
perturbation effect $\delta_d$ is applied to every cell,

\[
\widehat{x}_{i,d}
=
x_i+\delta_d,
\]
then the induced response is identical in every cellular state,

\[
r_{d,0}
=
r_{d,1}
=
\delta_d,
\]
and therefore

\[
\cos
\left(
r_{d,0},
r_{d,1}
\right)
=
1.
\]
An additive predictor may therefore preserve substantial heterogeneity in the untreated population
while generating no state-specific variation in the perturbation effect itself. This distinction
is why population heterogeneity and differential perturbation response are measured separately.

\paragraph{Latent-model reconstruction reference.}
For predictors that reconstruct the input before applying a latent perturbation shift, the measured
response depends on the reference used to define the perturbation effect. Let $P(C_k)$ denote the
model reconstruction of control state $C_k$, and let $W\delta_z$ denote the decoded latent
perturbation shift. Referencing the predicted population to the observed control cells gives

\[
r^{\,\mathrm{obs}}_k
=
\operatorname{mean}
\left(
P(C_k)
\right)
-
\operatorname{mean}
\left(
C_k
\right)
+
W\delta_z,
\]
whereas referencing it to the model's own reconstructed control gives

\[
r^{\,\mathrm{model}}_k
=
W\delta_z.
\]
The first quantity includes state-dependent reconstruction error; the second isolates the
perturbation effect represented by the model. Both references were therefore retained for latent
predictors. The additive-limit identity above applies to the model-referenced response, for which a
shared latent shift necessarily yields a cosine of 1.

\paragraph{Additional predictor diagnostics.}
We also evaluated deployed latent-arithmetic and non-additive transport-based generators under the
same state partition. These analyses test whether predicted populations retain state-specific
responses, not the relative quality of prediction architectures. Full implementations, acceptance
criteria and diagnostic results are reported in Supplementary Note~4 and Supplementary Table~4.


\subsection*{Recoverability of state-specific responses}

Differential responses can contribute to population-level retrieval only if the corresponding
cellular structure can be recovered from the observed response population. We therefore measured
recoverability by comparing unsupervised partitioning with a supervised estimate of how well the
same labels could be distinguished when their identities were known. This comparison separates
failure of an unsupervised method from absence of discriminative information in the measured
representation.

Two recovery problems carry this name in this paper. Both are balanced two-class problems scored
against a supervised ceiling and differ in what the two classes are, so results are not carried
between them. \emph{Response detectability} asks whether a drug's effect is legible in single cells at all:
its two classes are that drug's treated cells and the vehicle cells of the same context. It is the
quantity of Fig.~\ref{fig:5}f, and its supervised arm $A_{\mathrm{sup}}$ is the surviving explanator
of Fig.~\ref{fig:5}h, so that explanator measures whether the response is present in the cells and
not whether an unsupervised method finds it.
\emph{Response-identity recoverability} asks the harder question of whether cells can be assigned
to \emph{which} perturbation produced their response: its two classes are two drugs, or two drug
classes. It is the quantity of the constructed HDAC-versus-JAK two-state mixtures, of the
within-patient drug-pair splits in Fig.~\ref{fig:4}e, and of the Tahoe-100M drug-pair analysis in
Supplementary Note~3. Neither is the recovery of a cell state. Cell-cycle phase is itself called
from expression, so asking a probe to recover it from expression would be partly circular, and the
comparison this paper makes between supervised and unsupervised recovery is therefore posed on
externally determined labels throughout.

\paragraph{Controlled response mixtures.}
Controlled HDAC-versus-JAK mixtures were evaluated only as a positive control: source labels were
withheld from unsupervised procedures and used to estimate a supervised ceiling. The construction,
clustering procedures and scoring conventions are reported in Supplementary Note~1.

The supervised ceiling estimates the response-identity information present when labels are known;
unsupervised recovery measures how much of that information is extracted without labels. Full
cross-validation, drug-pair and state-separation sensitivity analyses are in Supplementary Note~1.


\fi
\subsection*{Natural heterogeneity and large-scale validation}

We next examined the information requirements for population-level retrieval in datasets in which
cellular heterogeneity was not created by mixing experimentally distinct source populations.
Patient-derived glioblastoma was used to test the requirements in naturally heterogeneous tissue,
and Tahoe-100M was used to assess the same quantities at substantially larger scale. The
glioblastoma cohort was not used as a retrieval benchmark; Tahoe-100M additionally served as the
substrate for the source-to-target intervention-identification benchmark described below.

\paragraph{Patient-derived glioblastoma.}
The ZhaoSims2021 dataset\cite{dataZhao} contains acute-slice cultures and surgical biopsies from ten glioblastoma
patients exposed to six compounds: Ana-12, Ipsenisib, RO4929097, Tazemetostat, etoposide and
panobinostat. Expression counts were normalized to $10^4$ per cell, log1p-transformed and reduced
to 2,000 highly variable genes. Because only one patient received more than two compounds, the
cohort does not support a sufficiently replicated leave-one-drug-out retrieval benchmark. We
therefore use it only to measure state-specific response, recoverability and the extent to which
population-level similarity differs from compartment-specific response ordering.

\paragraph{Cellular-compartment assignment.}
Cells were assigned to malignant and myeloid compartments using canonical marker scores; ambiguous
cells were excluded. Drug responses were defined relative to patient- and compartment-matched
controls. Assignment criteria, retained-cell counts and threshold-sensitivity analyses are reported
in Supplementary Note~3.

\paragraph{Differential response in natural tissue.}
State-specific perturbation responses were measured between malignant and myeloid compartments.
For patient $p$, drug $d$ and compartment $k$, the response vector was

\[
r_{p,d,k}
=
\operatorname{mean}
\left(
X^{\mathrm{treated}}_{p,d,k}
\right)
-
\operatorname{mean}
\left(
X^{\mathrm{control}}_{p,k}
\right).
\]
Only patient--drug combinations with at least 50 treated cells in both compartments were included.
Controls were matched within both patient and compartment. Differential response was summarized as

\[
R_{p,d}
=
\cos
\left(
r_{p,d,\mathrm{malignant}},
r_{p,d,\mathrm{myeloid}}
\right).
\]
This construction removes pre-existing differences between patients and cellular compartments
before comparing the directions of their drug-induced responses.

\paragraph{Recoverability in natural tissue.}
Response-identity recoverability was assessed from within-patient drug-pair splits by comparing
unsupervised partitions with a supervised ceiling. Patient-level bootstrap and method-pair
sensitivity analyses are reported in Supplementary Note~3.

\paragraph{Proxy for downstream decision relevance.}
The glioblastoma cohort does not contain enough drug coverage to test whether population-level
retrieval directly changes the preferred candidate. We therefore quantified a more limited
property: whether drug similarity at the whole-population level preserves the ordering observed
within a specific cellular compartment.

For each pair of drugs measured in the same patient, we first computed the cosine similarity
between their population-mean response signatures and the cosine similarity between their
malignant-compartment response vectors. Agreement between these two similarity structures was
summarized by Spearman correlation across drug pairs.

That whole-population comparison gives a Spearman correlation of $\PREMISEOVERLAP$, but the
population-mean signature is the equal-weight average of the two compartment responses, so the
malignant compartment contributes half of it and part of the agreement is arithmetic. We therefore
additionally used a disjoint comparison. For drug pair
$(a,b)$, myeloid-response similarity was defined as

\[
S^{\mathrm{myeloid}}_{a,b}
=
\cos
\left(
r_{a,\mathrm{myeloid}},
r_{b,\mathrm{myeloid}}
\right),
\]
and malignant-response similarity as

\[
S^{\mathrm{malignant}}_{a,b}
=
\cos
\left(
r_{a,\mathrm{malignant}},
r_{b,\mathrm{malignant}}
\right).
\]
We then computed

\[
\rho_{\mathrm{disjoint}}
=
\rho_{\mathrm{S}}
\left(
S^{\mathrm{myeloid}}_{a,b},
S^{\mathrm{malignant}}_{a,b}
\right),
\]
using only drug pairs available within the same patient. The two sides of this comparison contain
no shared cells and each compartment is referenced to its own matched control. This statistic is
interpreted as a proxy for the amount of candidate-ordering headroom available to population
structure, not as a direct measure of retrieval performance. No pass--fail threshold was applied.

\paragraph{Large-scale validation in Tahoe-100M.}
Tahoe-100M plate 3 was used to test the same requirements at larger scale. We evaluated two
independently defined states, released cell-cycle labels and a control-derived $k$-means partition,
and quantified state-specific response alignment, response-identity recoverability and agreement
of drug ordering across disjoint states. Dataset filtering and all state-definition details are
reported in Supplementary Note~3.

For each eligible drug--cell-line condition, the response of state $k$ was defined relative to its
own state-matched controls,

\[
r_{d,c,k}
=
\operatorname{mean}
\left(
X^{\mathrm{treated}}_{d,c,k}
\right)
-
\operatorname{mean}
\left(
X^{\mathrm{control}}_{c,k}
\right).
\]
Differential response was measured as the cosine between the two state-specific response vectors,

\[
R_{d,c}
=
\cos
\left(
r_{d,c,0},
r_{d,c,1}
\right).
\]
A condition was included only when at least 50 cells were available in each required treated and
control state.

Finally, a disjoint state-ordering statistic was used as a large-scale proxy for downstream
decision relevance. Within each cell line, cosine similarity was computed for every eligible pair
of drug responses separately in the two non-overlapping cellular states. Spearman correlation
between the two resulting drug--drug similarity matrices was then calculated,

\[
\rho_c
=
\rho_{\mathrm{S}}
\left(
S^{(0)}_{a,b},
S^{(1)}_{a,b}
\right).
\]
High correlation indicates that the two cellular states preserve similar relative ordering among
drug responses, whereas lower correlation leaves greater scope for state-specific information to
change candidate preference. As in the glioblastoma analysis, this quantity does not directly test
whether population-level retrieval improves a drug ranking.

These are intrinsic-state diagnostics and are reported separately from the intervention-identification
benchmark below, which uses the same plate. Detailed recoverability methods are in Supplementary
Note~3.


\subsection*{Source-to-target intervention identification on Tahoe-100M}

The retrieval experiments above rank candidates whose responses in the query context are observed.
An intervention task instead asks which compound will carry a starting population into a desired
one, so the candidate's effect in that context is unknown. The scope of that question here is
narrower than intervention design, and the naming follows the scope. The desired population is a
real treated population generated by one compound of the library, and the outcome is the rank that
compound receives, so the task measures whether population information helps identify the
intervention that produced an observed response. It does not measure therapeutic utility, and no
target state is specified independently of a compound that already reaches it. We defined this task on Tahoe-100M plate
3.

\paragraph{Task definition and eligibility.}
A context is one cell line. For a held-out context $c$ and a query drug $q$, the retriever receives
the vehicle (DMSO) cells of $c$ and an observed treated population $X_{q,c}$, and ranks all
\NCAND\ candidate compounds; it never sees $q$. A condition entered the grid if it had at least 100
treated and 200 vehicle cells, which admitted \NQUERIES\ eligible queries over \NCTX\ usable cell
lines from the 4,600 conditions present on the plate. Each query was evaluated at \NSEEDS\ fixed
seeds (13, 29, 47, 71, 101), giving \NQSEED\ query--seed pairs. Source and target populations were
drawn without replacement at up to 200 cells each, all available cells being taken where a
condition held fewer, and responses were referred to the source mean,
$P=\{x_i-\bar{x}_{\mathrm{source}}\}$.

\paragraph{Oracle arm.}
In the oracle arm every candidate's response in the held-out context is taken from the observed
data, which removes the forward model from the loop and bounds what any predictor could deliver.
Cell disjointness is enforced through the library preparation: the target population is drawn from
one sublibrary half and every candidate bank from the other, at up to 100 cells per candidate,
all cells in that half being taken where a candidate has fewer, without
replacement, so the generating drug is drawn from the same half as its 91 competitors and is not
penalized by a technical shift. This arm is an information ceiling and not a deployment result.

\paragraph{Prediction arm.}
In the prediction arm no treated cell of the held-out context enters any fit or any ranker.
Candidate responses are produced by a forward predictor fitted on the other \NFITCTX\ cell lines.
%% SIX, NOT FIVE. This paragraph named five predictors until 2026-09-04 and omitted
%% nearest_context_cells, while Fig. 6c drew and labelled all six that
%% analysis/phase2_transition/phase_b_loco.PREDICTORS actually runs. The omitted one is the only
%% predictor in the set that transfers real cells, so it is also the one the additivity argument
%% below most needs. The predictors are now grouped by that argument rather than listed flat.
Six predictors were evaluated. Three are additive: they emit one signature per drug and apply it
identically to every source cell, so the candidate population is the vehicle population rigidly
translated. These are an average effect $\delta_d$ pooled over training contexts; a nearest-context
transfer taking $\delta_d$ from the training context whose vehicle mean is closest in cosine; and a
linear-latent variant truncating the average effect to the principal components carrying 80\% of
training variance. Three are not additive. The \emph{state-conditioned} predictor estimates the
same average effect separately for each of the three cell-cycle states released with the plate, G1,
S and G2M, and applies it to each source cell according to its own state, which costs one extra
grouping and no new hyperparameter. Each per-state effect is averaged over the training contexts in which that state holds at least
20 cells; groups below the floor are dropped from the average rather than estimated, and 359 of
the 12,275 cell-line, drug and state groups fall below it, 294 of them in S phase. A drug left
with no usable estimate for a state would keep its pooled effect there and be additive for that
state alone; this did not occur for any of the $44\times92\times3$ signature slots, so the
predictor is non-additive over its whole output. The \emph{Gaussian optimal-transport} map
is a Bures map fitted in a 50-component training principal-component basis, from the
pooled within-context vehicle covariance to the drug's pooled within-condition treated covariance,
with the mean shifted by the average latent effect and applied to the held-out vehicle cells as a
latent edit so that the part of each cell outside the subspace is carried through; every cloud is
centred on its own condition mean before pooling, so a covariance here is within-condition cell
scatter. \emph{Nearest context, cells} transfers the neighbouring context's actual treated cells
for the drug, shifted by the difference of the two vehicle means; it is the cheapest predictor in
this set that can carry a non-additive response shape, and it is included because a set of purely
additive predictors would answer the question by construction rather than by measurement. Each
principal-component basis is fitted on up to 50,000 cells sampled from the training contexts, and
each context--drug fit uses up to 200 cells. The additive three generate no state-specific
differential response, which is an algebraic property of the model class and not a training
failure.

\paragraph{Scoring and evaluation.}
Candidates were scored by the three rules of Fig.~\ref{fig:1}a: a direction-only mean cosine, a
magnitude-aware mean $L_2$, and the population route, primarily the unbiased U-statistic energy
distance, with MMD and sliced Wasserstein as secondary variants. The secondary scores use settings
specific to this task: sliced Wasserstein over 200 random projections at 100 quantiles, and a single
RBF bandwidth fitted once by the median heuristic on the candidate library and held fixed across
queries, rather than the 128-projection and pooled multi-scale settings of the observed-population
arm. Mean reciprocal rank of the true
intervention is the primary outcome; Hit@1, Hit@5 and Hit@10 are secondary. Intervals are 2,000-fold
bootstrap intervals resampled over cell-line clusters, which are the independent units of this
design. A decision is counted as corrected when the population route places the true drug first and
the reference does not, and as broken in the reverse case.

%% ADDED 2026-09-04. Figures 5e, 5f and 5h had no Methods at all: the headroom null, the
%% detectability probe and the conditional regression were each explained only in the Results
%% sentence that cited them, which is not a specification a reader could re-run.
\subsection*{Headroom, response detectability and the conditional bottleneck analysis}

\paragraph{Headroom and its null.}
Headroom is defined per query as $H=1-\mathrm{RR}_{\mathrm{ref}}$, the room the direction-only mean
route leaves for any other scorer to improve on, and the gain is
$G=\mathrm{RR}_{\mathrm{pop}}-\mathrm{RR}_{\mathrm{ref}}$ on the same query. A Spearman correlation
between $H$ and $G$ is partly arithmetic, because the two share $\mathrm{RR}_{\mathrm{ref}}$ with
opposite signs and reciprocal rank is capped at 1. The null quantifies exactly that. Within each
cell line, the population route's reciprocal ranks are randomly permuted across queries, which
destroys any pairing between the two scorers while preserving both marginal distributions, the
$\mathrm{RR}\le1$ ceiling and the cell-line clustering; $G$ is then recomputed against the
unpermuted reference and the correlation retaken. Two hundred permutations at a fixed seed give the
interval reported in Fig.~\ref{fig:5}e.

\paragraph{Response detectability.}
For each held-out cell line $c$ and query drug $q$, equal numbers of cells treated with $q$ and
vehicle cells of $c$ form a balanced two-class problem, at 300 per arm or at the size of the
smaller pool where that is below 300, sampled at a fixed seed, so chance is 0.5 by construction. The supervised ceiling $A_{\mathrm{sup}}$ is the out-of-fold accuracy of a
five-fold cross-validated logistic probe ($C=1$) on those cells. The unsupervised value
$A_{\mathrm{unsup}}$ is $k$-means with $k=2$ on the same cells, scored by best-permutation accuracy,
so cluster-to-label matching is free and the unsupervised arm is not penalized for arbitrary label
order. The labels are external, being which well a cell came from, rather than an expression-derived
cell state. The supervised arm is a ceiling on the information present in the representation.

\paragraph{Conditional bottleneck regression.}
Regressing the gain $G$ directly on candidate explanators inherits the same arithmetic as the
headroom correlation. The specification reported in Fig.~\ref{fig:5}h therefore takes the
population route's own reciprocal rank as the outcome and includes the mean route's reciprocal rank
as a covariate, so that a coefficient answers a conditional question: given how well the mean route
did on this query, does this explanator predict how well the population route did. The covariate is
a control and is not plotted. The explanators are the interaction share, its cross-fitted
reproducibility $R_{\mathrm{int}}=\cos(I_A,I_B)$, response detectability $A_{\mathrm{sup}}$, and,
in a separate model reported for contrast, the naive response cosine $1-\cos(r_1,r_2)$. Estimation is
ordinary least squares with cluster-robust standard errors at the cell line, which is the
independent unit of the design, and all predictors are standardised to zero mean and unit variance
so that coefficients are comparable across explanators on one axis. Queries missing any explanator
are dropped listwise, which is why the reported $n$ of 3,618 is below the \NQUERIES\ of the
retrieval arm.

\subsection*{Cross-fitted drug-by-state interaction}

Testing whether a drug's effect depends on cell state requires a statistic that measures
interaction rather than effect size. Writing $r_s$ for the state-$s$ response, the difference
between mean treated and mean vehicle cells of that state, with states 1 and 2 the G1 and G2M
cell-cycle calls released with the plate, the quantity of interest is the interaction vector
$I=r_1-r_2$. Its squared norm cannot be used directly: with finite arms,
sampling noise alone makes $\lVert I\rVert$ large, and the same mechanism inflates
$1-\cos(r_1,r_2)$ for the weakest drugs.

Each of the four arms is therefore split into disjoint halves $A$ and $B$, the interaction is
estimated once in each, and the two estimates are multiplied:
\[
S_{\mathrm{int}}=\frac{\langle I_A,\,I_B\rangle}{p},
\qquad
S_{\mathrm{main}}=\frac{\langle M_A,\,M_B\rangle}{p},
\qquad
M=\tfrac{1}{2}(r_1+r_2),
\]
with $p$ the number of genes. Under the additive null $E[I]=0$ and the two halves are independent,
so $E[S_{\mathrm{int}}]=0$ exactly: noise makes $I_A$ and $I_B$ point in unrelated directions and
their inner product averages to zero, whereas a real state-dependent effect survives because both
halves estimate the same non-zero $I$. The vehicle arms are split as well; sharing one vehicle mean
between halves would place the same control noise in $I_A$ and $I_B$ and correlate them under the
null. Splits are library-preparation-disjoint wherever a condition spans at least two sublibraries,
and a plain random half of the cells otherwise.
Arms were capped at 400 cells and required at least 50, and ten A/B splits were averaged per query.

Two derived quantities are used. The \emph{interaction share},
$S_{\mathrm{int}}/(S_{\mathrm{int}}+S_{\mathrm{main}})$, is scale-free and has a cross-fitted
numerator and denominator, so neither term is noise-inflated; dividing by a plain response norm
would reintroduce exactly the single-sample denominator that sank the cosine statistic.
$R_{\mathrm{int}}=\cos(I_A,I_B)$ asks whether the interaction reproduces across halves. As a
positive control, synthetic conditions were built from the plate's own cells by taking the G1 cells
from one drug and the G2M cells from another, so interaction is present by construction in material
with the same states and arm sizes as the queries.

The cosine statistic $1-\cos(r_1,r_2)$ is retained only as a descriptive sensitivity analysis and
as the historical comparator in Fig.~\ref{fig:5}g. It is not used as a differential-response
measure anywhere that a conclusion depends on it. Where earlier analyses in this study used the
induced-response cosine to characterize predictors or tissue, that use is descriptive alignment
between two response directions and is labelled as such.


\subsection*{Statistical analysis}

Statistical units follow the experimental design rather than the number of cells or
query--candidate comparisons: cells estimate a perturbation-response population but are not
independent replicates of retrieval performance\cite{PseudoRep}.

The partial-observation arm departs from this rule: each query is treated as the unit of analysis although the
queries are not independent. Of the \NALL\ queries, 720 hold out one of 143 drugs and the remaining 45 pool the
unmeasured drugs of the partial-library split into a single batch query; the \NMOA\ leave-drug-out queries on
which mechanism recovery is scored come from 130 of those drugs. Each drug is scored across three shared cell
lines, several library fractions and up to five random seeds, so the \NALL\ queries rest on 420 distinct
drug-by-cell-line contexts. No cluster-level resampling was applied, so the $P$ values and intervals reported for
that arm are anticonservative with respect to the drug as a unit.

Paired comparisons between retrieval methods on the same queries used two-sided Wilcoxon
signed-rank tests; independent groups used Mann--Whitney $U$ tests; monotonic associations used
Spearman rank correlation. Across pre-specified strata, $P$ values were adjusted by the
Benjamini--Hochberg procedure\cite{BenjaminiHochberg}. Effect sizes are reported with every test; achieved power and the sample size required for
80\% power are computed for the one stratum where they bear on the conclusion and are reported in
Supplementary Note~1. Mechanism-of-action analyses were restricted to the 600 queries whose mechanism class
remained represented in the candidate library; undefined queries were excluded rather than encoded
as zero.

For the GDSC2 functional analysis, Spearman correlations were computed per query across its
candidate drugs and summarized within each SciPlex3 cell line, which is the analysis unit;
query-level paired tests are reported as nominal. The protein evaluator-shape experiment was
summarized within the Control, IFN$\gamma$ and Co-culture conditions, its three higher-level units.
For HIR-Bench, cross-validation was grouped by the label-determining
$(\alpha,\mathrm{conflict})$ pair, giving 28 effective parameter cells, and performance was
summarized by pooled out-of-fold AUC. For the patient-derived glioblastoma analysis, uncertainty in
the recoverability gap was estimated by cluster bootstrap at the patient level with 5,000
resamples, differences being taken within each drug-pair split before aggregation; four patients
contributed evaluable splits. Intervention-retrieval intervals are 2,000-fold bootstrap resamples
over cell-line clusters. Tahoe-100M summaries are taken across eligible drug--cell-line conditions
or within-cell-line drug pairs according to the statistic, and both cellular-state definitions,
released cell-cycle state and the control-derived $k$-means partition, were analysed in parallel.

Analyses were implemented in Python. Single-cell data handling used Scanpy\cite{Scanpy}; the
$k$-means, Gaussian-mixture, principal-component and logistic-regression procedures named above
are the scikit-learn implementations\cite{ScikitLearn}.


\section*{Data availability}

\sloppy

All datasets analysed in this study are publicly available.

SciPlex3, Frangieh Perturb-CITE-seq and ZhaoSims2021 were obtained in harmonized form
from scPerturb\cite{ref3}, Zenodo record 13350497
(\href{https://doi.org/10.5281/zenodo.13350497}{doi:10.5281/zenodo.13350497},
release 1.4). The exact harmonized files used in the analyses were
\texttt{SrivatsanTrapnell2020\_sciplex3.h5ad} for SciPlex3,
\texttt{FrangiehIzar2021\_RNA.h5ad} and
\texttt{FrangiehIzar2021\_protein.h5ad} for Frangieh Perturb-CITE-seq,
and \texttt{ZhaoSims2021.h5ad} for the patient-derived glioblastoma analysis.
The primary SciPlex3 data are available from the Gene Expression Omnibus under
accession GSE139944\cite{dataSciPlex3}; the original Frangieh and ZhaoSims2021
datasets are described in the corresponding publications\cite{dataFrangieh,dataZhao}.
The CD34$+$ haematopoietic progenitor dataset is available from the Gene Expression
Omnibus under accession GSE306429\cite{dataCD34}.

Tahoe-100M\cite{Tahoe100M} was released by Vevo Therapeutics and the Arc Institute
under CC0 and is distributed through the Arc Virtual Cell Atlas and Hugging Face
(\texttt{tahoebio/Tahoe-100M}). This study uses plate 3 only, together with the
released \texttt{obs\_metadata} and \texttt{sample\_metadata} tables. Analyses use
the released preprocessed expression matrix comprising 4,158,278 cells and 2,304
highly variable genes in log1p-transformed space.

The external functional evaluation used GDSC2 fitted dose--response data from the
Genomics of Drug Sensitivity in Cancer, release 8.5
(\texttt{GDSC2\_fitted\_dose\_response\_27Oct23.xlsx}), available from
\url{https://www.cancerrxgene.org}. The compound-matching table used to align the
SciPlex3 and GDSC2 drug panels (\texttt{drug\_match\_table.csv}) is provided with the
analysis code.

Processed analysis inputs and the scripts used to construct them from the source
datasets are provided in the accompanying repository. Source-data tables underlying
each figure panel are provided with the corresponding analysis outputs.


\section*{Code availability}

All code used for data processing, retrieval, evaluation and statistical analysis is
available at \url{https://github.com/Boom5426/PopRetrieve}. 

\section*{Acknowledgements}
\TODO{Acknowledgements. Funding sources and grant numbers.}

\section*{Author contributions}
\TODO{Author contributions, CRediT taxonomy: conceptualization, methodology, software,
formal analysis, investigation, data curation, writing (original draft), writing (review and
editing), visualization, supervision, funding acquisition.}

\section*{Competing interests}
\TODO{Competing interests statement. If none: ``The authors declare no competing interests.''}

\section*{Additional information}
\noindent Supplementary Information is available for this paper: Supplementary Methods, Supplementary Notes 1--4 and
Supplementary Tables 1--4.\\
Correspondence and requests for materials should be addressed to \TODO{Name}.

\bibliographystyle{unsrtnat}
\bibliography{references}

\end{document}
